\documentclass[11pt]{article}

\usepackage[margin=1in]{geometry}
\usepackage{amsmath, amssymb}
\usepackage{booktabs}
\usepackage{array}
\usepackage{microtype}
\usepackage{xcolor}
\usepackage{titlesec}
\usepackage{enumitem}
\usepackage[hidelinks]{hyperref}
\usepackage{caption}
\usepackage{float}

\captionsetup{font=small, labelfont=bf}
\titleformat{\section}{\large\bfseries}{\thesection}{0.6em}{}
\titleformat{\subsection}{\normalsize\bfseries}{\thesubsection}{0.6em}{}
\setlist{itemsep=2pt, topsep=3pt}

\newcommand{\real}{\mathbb{R}}
\newcommand{\arm}[1]{\textsf{#1}}

\title{\vspace{-2.5em}\bfseries Multi-Channel Hebbian Plasticity in Multilayer BDH:\\
A Learned Partition of Working Memory\\ on a Synthetic Binding Task with Local Cues\\[0.3em]
\large\normalfont A window gate, slow memory and a channel split aimed at the task's key positions\\ find an eight-way partition of the grouped layout without stream labels or restarts,\\ on two machines with disjoint seeds, and $17/20$ in a replication from the text}
\author{A technical report \quad\textbullet\quad Revision 9 (draft)}
\date{Project repository: \texttt{Ynkling/clankers}}

\begin{document}
\maketitle
\vspace{-1.5em}

\begin{abstract}
BDH's Hebbian working memory is split into $k$ channels, and a per-token gate distributes each token's write over the channels and blends its read from them. On a synthetic task, $S$ streams bind the same keys to different values, cued by a stream token; one channel cannot hold the conflicting bindings at four or eight streams, and a hand-set gate, one channel per stream, can. The question is whether a gate trained only on the task loss finds that partition. Revision~8 answered it for the task's grouped layout, where the cue is within three tokens of every write: a \emph{window gate} (a width-3 causal convolution of the embeddings), slow memory for the first 2400 updates and a plateau-triggered \emph{channel split} that copies the busiest channel's gate row onto the idlest bind eight streams, the streams on disjoint sets of channels, in $38/40$ runs of a pre-registered test on two machines with disjoint seeds, against $0/40$ for the recurrent-gate recipe (38 vs 0 discordant pairs) and $15/40$ without the split (23 vs 0). The recipe uses no stream labels, no restarts and no hinge; the split reads the task's key positions.

This revision reports what has been done to that result since. An independent referee session recomputed the numbers of the report's tables from the records (140 of 156 reproduce; the seven that disagreed were statements, corrected in Revision~8.1), audited the pre-registration of every claim (none edited after a run), compiled the project's claim registry (53 pre-registered tests per machine, of which the headline survives any multiplicity correction), and found statements in Revision~8 that were wrong when committed; they are corrected and listed. A clean-room replication from the report alone, by a session that never saw the code, bound eight streams $17/20$ against $7/20$ without the split (10 vs 0): the model it built agrees with ours to $7\times10^{-9}$ in the gates and $7\times10^{-8}$ in the logits, the 44 details it had to guess are stated in an appendix, and one of them exposed a sampling weakness in the split's own probe. On the main line, the second claim of Revision~8 is now pre-registered and, on one machine, shown as that machine's secondary claims (the other machine's half has not run, and the pooled readings wait for it): the copy with its reset binds eight streams $10/10$ against the reset alone $3/10$ (7 vs 0), the reset alone is within the bound of no trigger ($3/10$ against $5/10$), and one delta-rule channel binds $0/10$ where the recipe binds $10/10$. Exploratory screens, which are not claims, answered three questions the earlier revision left open, each on ten to twenty paired seeds. The delta-rule write, with its write strength fixed rather than learned, is not shown worse than the Hebbian write at two, four and eight streams ($10/10$, $10/10$ and $9/10$ against $10/10$) and binds sooner at two and eight, so the memory rule is open again and is the next main-line test, specified and reviewed. On Hebbian memory the layer whose routing matters is the first: perfect gates at the later layers alone bind $0/10$, and a register that overwrites at the context token, read by a gate at all three layers, carries a cue up to six tokens away $20/20$ when nudged by labels, is found label-free on $4/20$ seeds, fourteen of the sixteen failures on each memory being key splits, formed in the records before any write at the cue, and is found $20/20$ on fresh seeds when a hinge on the gate's key-dependence holds the key split off for the first 2400 updates, a training-time prior that reads the task's key positions and identities, as the split reads key positions. And under a dense next-token loss on interleaved sources, a single channel learns to carry the source by itself. Two-stream key splits resist repair and are prevented by the hinge's key term on a screen ($20/20$ against $15/20$), as the register's are; a replication's probe can be wrong about a plateau; the title is scoped to what was shown.
\end{abstract}

\paragraph{Revision~8 (7--10 October 2026).} Revision~8 added Phase~VI: two pre-registered tests on two machines with disjoint seeds, an audit of the project's records (\texttt{docs/audit\_2026-10.md}), five exploratory batches on the side branch, three exploratory sessions on their own branches and a reading of twelve papers (\texttt{docs/reading/}). Its first draft was reviewed by an independent referee session, which found statements that were wrong when committed; Revision~8.1 corrected them and nothing else (the list is in Section~\ref{sec:referee}).

\paragraph{Revision~9 (draft).}\sloppy This revision adds, on the main line, a third pre-registered test (\texttt{test\_split\_copy}, complete on X; L's half has not run) and a fourth (\texttt{test\_delta\_recipe}), specified, reviewed by the referee, amended once and written, its test file checked against the amended spec before any run; the referee's report, recomputation and claim registry, merged into the repository (\texttt{docs/review\_2026-10/}); an independent clean-room replication of the main result (branch \texttt{claude/replicate-R}); the exploratory sessions' second and third rounds (F's S68--S81, G's S55--S83, H's S61--S86, I's S74 and a GPU package); and L's result files for Phases~III--IV, which were missing from the repository. Two machines ran every main-line test: \textbf{X}, a cloud container on Intel Xeon hosts at 2.10 and 2.80\,GHz, and \textbf{L}, an Intel i7-12650H. From the Muon-recipe test on, X and L run disjoint seeds and the pooled claim is primary. The exploratory sessions are \textbf{E} (branch \texttt{claude/outside-ideas}, batches 1--19), \textbf{F}, \textbf{G}, \textbf{H} and \textbf{I} (branches \texttt{claude/explore-F/G/H/I}), the replication \textbf{R} (\texttt{claude/replicate-R}) and the referee \textbf{V} (\texttt{claude/review-V}); F, G, H, I and V ran on 2.10\,GHz Xeons (R's records carry no CPU model), E's batches mostly on 2.10\,GHz hosts with 60 of batch~16's 165 runs and batch~19 on 2.80\,GHz. Every verdict was fixed before its run and printed mechanically; the referee's recomputation (Section~\ref{sec:referee}) re-derived every pre-registered reading of the sessions' later reports from their records and checked that the earlier ones were committed before their first run. The referee also checked this revision's first draft against the records, and the statements it found wrong are corrected here and listed in Section~\ref{sec:referee}. The revision is a draft: L's half of \texttt{test\_split\_copy} has not run, E's batch~19 (S66--S67), F's S81, G's S83 and H's S86 are running and R's Phase~4 has not started, and everything remains one small model on synthetic tasks. The title has been scoped to what was shown.

\begin{table}[H]
\centering
\footnotesize
\setlength{\tabcolsep}{4pt}
\begin{tabular}{@{}l>{\raggedright\arraybackslash}p{0.27\linewidth}>{\raggedright\arraybackslash}p{0.36\linewidth}r@{}}
\toprule
test, batch or session & question & verdict & \S\\
\midrule
\texttt{muon\_recipe} (X, L, pooled) & the early hinge window under Muon; Muon's speed at four streams and eight keys & M1 SHOWN pooled; M2, M3 NOT SHOWN pooled; reading not met & \ref{sec:muon}\\
\texttt{window\_gate} (X, L, pooled) & the window gate with slow memory, no hinge, at 2, 4 and 8 streams; the split at 8 & G1, G2, G4 SHOWN pooled; G3 HOLDS; ``eight streams bind without labels or restarts'' & \ref{sec:main}\\
\texttt{split\_copy} (X; L to run) & the copy against the reset; one delta channel at eight streams & on X (secondary): C1, C3 SHOWN, C2 HOLDS; pooled claims and readings wait for L & \ref{sec:split}, \ref{sec:memory}\\
\texttt{delta\_recipe} (specified) & the delta write against the Hebbian write at eight streams, and the split on each & spec reviewed by the referee and amended; the test written and checked against the amended spec; not run & \ref{sec:memory}\\
\midrule
batches 15--18 (E) & the recurrent gate's cap; the window gate; the split; the reset control; key splits & as in Revision~8 & \ref{sec:window}, \ref{sec:split}\\
batch 19 (E, running) & width and depth at four streams; eight streams at the larger sizes & scale-free to $N = 1024$, $D = 128$, 4 layers at four streams & \ref{sec:scale}\\
session F & delta-rule channels; the premise test; forget gates; the lite task; the probe; the delta recipe & one delta channel $2/10$, $0/10$ at four and eight streams; the delta recipe with fixed $\beta$ keeps 2, 4 and 8 streams & \ref{sec:memory}\\
session G & the cue in the stack; the register; which layers must be routed & the cue is absent; layer~1 is the block on Hebbian memory; the register carries the cue $20/20$ nudged, $4/20$ label-free, $20/20$ with the key hinge & \ref{sec:far}\\
session H & anti-collapse devices; the copy without reset; the split's constants; two streams & neither device carried to eight streams replaces the split; the copy alone suffices (a one-pair bound); the cap matters; the flat-start guard moves the splits to higher probe accuracies without changing the count; two-stream key splits resist the copy and are prevented by the key hinge, not by a slow gate & \ref{sec:split}, \ref{sec:keysplits}\\
session I & a dense next-token loss on interleaved sources & one channel learns to carry the source; four sources need a GPU & \ref{sec:lm}\\
session R & a clean-room replication of \texttt{window\_gate}'s eight-stream result & $17/20$ against $7/20$ (10 vs 0): partially replicated by its own bar, consistent with $38/40$ & \ref{sec:replication}\\
session V & referee: recomputation, pre-registration audit, claim registry, literature, spec and draft reviews & 140 of 156 numbers reproduce, the rest statements; 53 tests per machine; Revision~8's and this draft's errors; the probe's imbalance & \ref{sec:referee}\\
\bottomrule
\end{tabular}
\caption{Phase~VI. Main-line tests are files \texttt{test\_<name>.py} whose docstring records the pre-registered design and, after the run, the result on each machine and pooled; from \texttt{test\_window\_gate} on, the specification is committed to \texttt{specs/} before the test is written, and from \texttt{test\_delta\_recipe} on it is reviewed by the referee session before any run. Exploratory batches and sessions write only new files on their own branches and label every output ``exploratory, not a result''. Appendix~\ref{app:prov} lists the commits.}
\label{tab:index}
\end{table}

\section{Background}
\label{sec:background}

BDH \cite{bdh} stores working memory in synaptic state updated by a Hebbian rule. In its GPU form the operation is linear attention: the score between a query position $t$ and a past position $s$ is a content match $x_t \cdot x_s$ between sparse positive neuron activations, weighted by a positional operator (scalar decay or RoPE). The public implementation, \texttt{bdh.py}, ties queries to keys ($Q = K$), shares one layer's weights across depth, and LayerNorms the attention output. Its state update, in the paper's notation, is $\rho_{t} = \rho_{t-1} + \mathrm{LN}(E y_{t})\, x_{t}^{\top} U$: a pure additive write.

The extension gives every synapse $k$ channels and every token a gate $g_t \in \Delta^{k-1}$ that distributes its write over the channels and blends its read from them. With one distribution for both, the effective score becomes
\begin{equation}
\underbrace{(x_t \cdot x_s)}_{\text{content match}}\;\times\;\underbrace{(g_t \cdot g_s)}_{\text{context-gate match}}\;\times\;\mathrm{decay}^{\,t-s},
\label{eq:score}
\end{equation}
strictly causal ($s < t$), so tokens sharing a context can read each other's memory while tokens in different contexts are isolated. The score has the three-factor form that Triadic Linear Attention \cite{tla} describes (cited for its form; the paper is not among the twelve read in full), with $g$ as the second key; there the second key is a learned projection of the token, here it is a routing distribution inferred from context, and the subject is its label-free discovery. With $k = 1$ the model is single-channel BDH. One gate $g_t$ is computed per token from the raw token embedding $v_t$ (not LayerNormed, not convolved, not the residual stream) and applied at every layer. The hand-set \emph{perfect gate} sends every token of stream $s$ to channel $s$; it is used only as a validity ceiling.

\paragraph{Two learned gates.} The \emph{recurrent gate} of Phases~II--V,
\begin{equation}
h_t = \tanh\bigl(W_{\mathrm{in}}\,v_t + W_h\,h_{t-1}\bigr), \qquad g_t = \mathrm{softmax}(W_g\,h_t),
\label{eq:gate}
\end{equation}
with 32 units in $h_t$ and $W_{\mathrm{in}}, W_h, W_g \sim \mathcal{N}(0, 0.1^2)$ (so that $\rho(W_h) \approx 0.6$ and $\sigma_{\max}(W_h) \approx 1.08$ at initialization). The \emph{window gate} of this revision,
\begin{equation}
h_t = \tanh\Bigl(W_{\mathrm{in}} \sum_{j<3} w_j \odot v_{t-j}\Bigr), \qquad g_t = \mathrm{softmax}(W_g\,h_t),
\label{eq:window}
\end{equation}
a depthwise causal convolution of width 3 over the embeddings (weights drawn uniformly in $\pm 1/\sqrt{3}$, no bias), with no recurrence. It sees the current token and the two before it, which in the grouped layout is enough to see each key's stream token at the key, the value and the query. It cannot see anything further back, and it has no notion of position.

\paragraph{The uniform gate is a stationary point.} Writing each gate as a deviation from uniform, $g_t = \tfrac{1}{k}\mathbf{1} + \delta_t$, gives $g_t \cdot g_s = \tfrac{1}{k} + \delta_t \cdot \delta_s$. The constant is cancelled by the LayerNorm after attention, so routing enters only through products of deviations, and at the uniform gate the task loss exerts no first-order pull toward any routing. This is why discovery is a race between the gate's first commitment and whatever the memory learns first, and why the recipes below act in the first few thousand updates.

\paragraph{A short causal convolution.} Most arms include the short convolution that Mamba and most linear-attention models place before their token mixer \cite{mamba}: a depthwise causal filter of width~4 at every layer, initialized to the identity, feeding the queries and the values. Without it, the perfect gate of this model ($N = 256$) binds the two-stream configuration on every seed recorded ($4/4$ in the main line's \arm{CEIL\_A}, $14/14$ in session~F). The first draft attributed four failures to the missing convolution; those runs used a smaller model (\texttt{test\_multilayer\_binding}'s $N = 64$ ceiling, 7{,}552 parameters against 25{,}984), as the referee's recount found. Whether the convolution carries part of the binding is open.

\section{Earlier phases in brief}
\label{sec:prior}

\paragraph{Phases I--IV (Revisions~1--6).} A gate that was a function of token identity could not express the routing (Phase~I); a recurrent gate could (Phase~II). On the binding task of Section~\ref{sec:task}, two channels with the perfect gate bind two streams where no single-channel model binds, and a label-free learned gate finds the routing on about half its runs (Phase~III). A restart rule made two streams reliable; four streams with four channels mostly failed by merging streams into a shared channel, and spare channels removed most merges (Phase~IV).

\paragraph{Phase~V (Revision~7).} An early-training recipe replaced restarts at two streams and at four with sixteen channels: for the first 2400 updates the memory trains at a tenth of the gate's rate (\textbf{SLOW}), and a hinge penalizes the gate when its channel choice at key positions is explained by position (\textbf{HINGE}; \textbf{WINDOW} switches the hinge off after update 2400). Eight streams failed in every arm (none of more than 100 runs), and screens located the failure in the recurrent gate: its state carries the stream at initialization and loses it within the first 600--1200 updates, as its recurrent term grows and, in seven of ten runs, the spectral radius of $W_h$ rises above one. Revision~7 ended with a cap on that recurrence being screened.

\section{The tasks, the outcomes and the recipes}
\label{sec:task}

\paragraph{The binding task.} There are $S$ streams, each with a context token $\mathrm{CTX}_s$. For each of $P$ keys, $S$ distinct values are drawn from 16 value tokens, so every key is bound to a different value in every stream. In the \emph{grouped} layout the body is $S P$ triples $[\mathrm{CTX}_s, \mathrm{KEY}_i, \mathrm{VAL}_{i,s}]$, grouped by key, keys in random order and stream order random within each key's group. One query $[\mathrm{CTX}_q, \mathrm{KEY}_q, ?]$ ends the sequence; the model predicts the value from the query's KEY position, where $\mathrm{CTX}_q$ is one token back. The write that must carry the stream is at the VAL position, where CTX is two tokens back. Chance among values is $1/16$; a model that binds keys to values but cannot tell streams apart scores $1/S$. The \emph{header} layout gives each stream one block $[\mathrm{CTX}_s, \mathrm{KEY}, \mathrm{VAL}, \mathrm{KEY}, \mathrm{VAL}, \ldots]$ with the stream token once at the start, so most writes are far from their cue; at a fixed block length, position can route it. The \emph{randomised header layout} (``RandHeaderTask-lite'', session F's S70; the distant-cue task of Sections~\ref{sec:far} and \ref{sec:replication}) has $S = 2$ streams of $P = 4$ keys, each stream's keys in two blocks whose lengths are a uniformly drawn composition of 4 into parts of 1--3 pairs, the four blocks in uniformly random order, one CTX token at the start of each block, every (stream, key) once, and one query; $L = 23$, and a token is one to six tokens from its block's cue, the write at a value position two to six. Both perfect gates bind it by update 1200 on five of six seeds. The model has three layers with weights shared across depth, decay $0.95$ (half-life about 14 tokens), $N = 256$, embedding width 32 and one head, and trains at batch 32 for a budget set per test (24{,}000, 28{,}800 or 43{,}200 updates; 48{,}000 in some of session~F's screens), with evaluation every 1200 updates on 2048 held-out queries; a run stops early after three consecutive evaluations at or above $0.95$. Session I's dense-loss task is defined in Section~\ref{sec:lm}.

\paragraph{Outcomes.}
\begin{itemize}
\item \textbf{Bound:} an evaluation reaches $0.95$ and every later one stays there; the \emph{transition} is the first such step. A run that reaches its budget counts as bound on a single final evaluation; this has happened in 14 of 1869 bound runs in \texttt{results/} at Revision~8's commit (the referee's recount).
\item \textbf{Discovered} ($k = S = 2$): bound, with the streams' read gates separated at value positions (final VAL cosine below $0.5$).
\item \textbf{Bound routed} ($k \geq S > 2$): bound, every stream has its own channel (a one-to-one \emph{stream-to-channel map}, measured at value positions, where the writes are), and every stream's accuracy is at least $0.9$. The map takes each stream's largest-mass channel, from a 256-sequence probe, ties resolved on unrounded values; in the 38 headline runs the streams' channel sets are disjoint (largest pairwise overlap $0.047$) while in 18 of them some stream's largest channel holds under $0.9$ of its gate mass. On the randomised header layout session~F's \textbf{ROUTED*@end} (routing margin $\geq 0.9$ and $\eta^2$ by stream at key positions $> 0.9$ at the end) is used, and the readings of session~G's S72, S73 and S82 (Section~\ref{sec:far}) are on \textbf{ROUTED-BOUND}, bound and ROUTED*@end, while S71's and S56's are on bound and on G's own bound routed; the name DISCOVERED is reserved for the outcome above, since one screen's reading named ``discovered'' was fixed on ROUTED-BOUND and would have come out differently on DISCOVERED (Section~\ref{sec:far}).
\item \textbf{Failure classes:} POSITION or KEY when the gate's channel choice at key positions is explained ($\eta^2 \geq 0.5$) by position or by key identity; STREAM-PARTIAL; OTHER. For $S > 2$: MERGED ($j$ share) when $j$ streams share a channel, non-stream when the gate does not split by stream, \emph{collapsed} when final accuracy is below $0.15$. The classifier of record (v1) applies the margin rule before the key rule and labelled one batch-18 run at $\eta^2 = 0.500004$ by key STREAM-PARTIAL (two more sit at $0.49986$ and $0.499985$ and keep that label under both); \texttt{fail\_class\_v2} (commit \texttt{ba8901a}) applies the key rule first and is used by every test from \texttt{test\_split\_copy} on; recorded verdicts keep the rule they were read under, and the two classifiers disagree only where $\eta^2$ by key $\geq 0.5$ and the margin rule fired (six of 29 recorded unbound two-stream runs of three earlier screens, four of S78's seventeen, and none at eight streams in S79).
\end{itemize}

\paragraph{Recipes.}
\begin{itemize}
\item \textbf{SLOW:} one Adam optimizer with two groups. The gate's parameters train at $10^{-3}$ throughout; every other parameter, embedding and convolution included, at $10^{-4}$ for updates 1--2400 and $10^{-3}$ after. The perfect gate (the validity arm) trains at a single rate, $10^{-3}$, without SLOW.
\item \textbf{HINGE, WINDOW:} Revision~7's position hinge, on the recurrent gate; kept here only as the comparison.
\item \textbf{SPLIT (KEYMASS):} at checks every 2400 updates from 4800, if the probe's accuracy is below $0.95$ and rose by less than $0.02$ since the previous check (the probe is also read at 2400 as the first reference; the last check is at the budget minus 2400), copy the gate row $W_g[c^*]$ of the channel with the largest mean read-gate mass at the body's key positions onto the row $W_g[c_0]$ of the channel with the smallest; at most three splits per run, at least 4800 updates apart. \texttt{test\_window\_gate}'s form adds noise of $0.1$ times the row's standard deviation to both rows and zeroes $W_g$'s Adam state, step counter included; session~H's \emph{exact copy} does neither, and is the form of the screens of Sections~\ref{sec:split} and \ref{sec:memory} and of the next main-line test. The probe is one fixed set of 64 sequences, the same for every seed and arm (\texttt{Generator(12345)}), scored on its queries only; its queries fall 5, 5, 6, 8, 10, 12, 5 and 13 on streams 0--7, so it is unbalanced in the same way for every run (Section~\ref{sec:replication}). The split uses no stream labels; it needs to know where the key positions are.
\item \textbf{MUON:} Muon at learning rate $0.005$ on every 2-D weight (momentum $0.95$, Nesterov), with Adam on the rest, and SLOW's schedule.
\item \textbf{The delta write (session~F):} each channel's state follows Gated DeltaNet's rule, $S^c_t = S^c_{t-1}\bigl(\alpha^c_t(I - g^c_t \beta_t \hat{x}_t \hat{x}_t^\top)\bigr) + g^c_t \beta_t v_t \hat{x}_t^\top$, with unit keys $\hat{x}_t$, the decay fixed at $0.95$, and the write strength $\beta$ \emph{fixed at 1}; at $\beta = 0$ its logits agree with the Hebbian model's to $2\times10^{-7}$, and with the new path disabled it reproduces a recorded run bit for bit through 2400 updates. A learned $\beta$ collapsed toward ``do not write'' at value positions in two screens and rose above its initial $0.5$ in two others (Section~\ref{sec:memory}).
\item \textbf{The register gate (session~G):}\sloppy a 16-dimensional register $R_t = R_{t-1}(1 - \beta^R_t) + \beta^R_t W_R v_t$, $\beta^R_t = \sigma(w_R \cdot v_t + b_R)$ with $b_R$ initialised to logit $0.1$, written by every token and read causally (the gate at $t$ reads $R_{t-1}$); one gate shared by all three layers, $g^{(\ell)}_t = \mathrm{softmax}(W_g \tanh(A\,[\mathrm{LN}(r^{(\ell)}_t);\, u_t;\, R_{t-1}]))$ over the layer's residual input, the window gate's input $u_t$ and the register, with $A$ initialised so that the gate equals the window gate at step 0 up to a $\mathcal{N}(0, 0.01^2)$ perturbation of its register and residual blocks, and a stop-gradient from the gate into the stack. ``REG3'' is this gate at all three layers; session~G's first design (S56) had a register gate at layers 2--3 only, reading the raw embedding rather than the window, with the window gate at layer~1.
\end{itemize}

\paragraph{Statistics.} Paired comparisons in the main-line tests use exact one-sided McNemar tests on the discordant pairs (``$b$ vs $c$''), unpaired ones Fisher's exact test; E's exploratory batches report two-sided McNemar $p$-values and sessions F, G and H one-sided ones, as marked where quoted; bands are RELIABLE at $\geq 90\%$ and MAJORITY at $\geq 50\%$ of runs, with Wilson 95\% intervals. From the Muon-recipe test on, X and L run disjoint seeds (X 400--599, L 1400--1599), the pooled claim over both machines' pairs is primary and each machine's is secondary. The referee's claim registry (\texttt{docs/review\_2026-10/claims\_registry.md}) counts 53 pre-registered hypothesis tests per machine over Phases~III--VI (52 one-sided; 25 Fisher, 25 McNemar, 2 Mann--Whitney, 1 sign test), plus six pooled tests in Phase~VI, besides 3 bounds without a $p$-value, 41 bands, 29 validity gates and 37 readings; every test was run at $\alpha = 0.05$ with no adjustment, which controls the per-comparison rate only, so about 2.7 false positives per machine are expected if every null were true. The conclusions are assessed under Benjamini--Hochberg at $0.05$ after the fact: 19 of X's 23 positive verdicts and 27 of L's 32 survive, and the four on X (\texttt{stream\_recipe} R3, \texttt{recipe\_scope} Q2, \texttt{slow\_start} S0, \texttt{muon\_recipe} M2) and five on L (\texttt{short\_conv} K2, \texttt{stream\_channels} MG, \texttt{slow\_start} HA, H0S and S0) that depend on the unadjusted level carry no conclusion below; the main result (G2, G4; the pooled G1 and M1) survives Holm on X, on L and pooled. The conclusions that matter are those replicated on both machines or pooled over disjoint seeds.
\section{The window gate}
\label{sec:window}

\paragraph{History.} The window gate was the first screen of the first exploratory batch (S1, batch~1): on seeds 160--179 with plain Adam and no recipe it discovered $15/20$ against the recurrent gate's $7/20$ (12 vs 4, two-sided $p = 0.077$: ``promising'' by the screen's rule) and bound sooner on average (mean transition 2720 against 3429; equal medians); its five failures were three key splits and two stream-partial runs under the main-line classifier, four key splits under the key-first rule (it cannot position-split, having no position). Batch~2 ran it on the header layout, where it bound $2/10$, both by key splits, and it was dropped: a gate that cannot see a cue three tokens back seemed the wrong direction for a language model. It was never combined with slow memory and never run at four or eight streams. The October audit of the project's records (\texttt{docs/audit\_2026-10.md}, Section~3.1) noticed that every obstacle of Phase~V was a property of the recurrent parameterization: position splits are the recurrence counting triples, and the eight-stream loss of the stream is the recurrence becoming expansive. A window gate can do neither. The audit fixed a prediction before any run: with slow memory, at two streams $\geq 36/40$ with key failures gone; at eight streams, binding on some seeds where every recurrent arm bound none.

\paragraph{Batch~16 (E).} The window gate with SLOW and no hinge (\arm{LOCAL3\_SLOW}): at two streams $31/40$ on seeds 160--199 (STREAM-PARTIAL 5, KEY 4, no POSITION), inconclusive against the hinge recipe's $34/40$ (5 vs 8). At four streams with sixteen channels, $20/20$ bound routed against X's recorded \arm{HINGE4k16} $15/20$ (5 vs 0, two-sided $p = 0.0625$). At eight streams, $3/10$ under Adam and $3/10$ under Muon, where every recurrent arm in batches 13--16 had bound $0/10$; the seven unbound Adam runs ended with two or three streams sharing a channel at accuracy $0.73$--$0.88$. The window gate's state decodes the stream at both key and value positions ($1.00$ at four streams at update 4800, the two runs still training at 9600 reading $0.70$--$0.72$; at eight, $0.93$--$1.00$ under Muon and, under Adam, $0.55$--$0.67$ at updates 2400--4800 recovering to $0.87$--$0.88$ by 9600), where the recurrent gate's reads $0.14$--$0.28$. So the audit's first prediction was partly met (no position failures; key failures reduced, not gone, and $31/40$ short of 36), and its second was met.

\paragraph{Batch~17 (E).} The failures at eight streams were merges, so batch~17 added the targeted split of Section~\ref{sec:task} (S37's KEYMASS rule, at most three splits). \arm{W\_SPLIT} bound routed $9/10$ against $3/10$ for the same seeds without it (6 vs 0, two-sided $p = 0.031$); every one of its 19 splits was on target by the labels (the busy channel held two to five streams, the idle one none), and in six of nine bound runs the transition came at the evaluation right after the last split. Decay $0.98$ alone bound $1/10$ (it settled the merges early) and a longer budget alone $4/10$: the splits, not the decay or the time, do the work. The screen handed \arm{W\_SPLIT} back for a pre-registered test.

\section{The main result: \texttt{test\_window\_gate}}
\label{sec:main}

The test's specification was committed before the test was written (\texttt{specs/test\_window\_gate.md}) and amended, before any Part~D run, to add the split. Both machines ran it on their own seeds (X 500--559, L 1500--1559). Part~A ($S = 2$, $P = 4$, $k = 2$, no convolution, 24{,}000 updates, 40 seeds per machine): the window gate with one learning rate (\arm{WIN3}), with slow memory (\arm{WIN3\_SLOW}), and the recurrent-gate recipe (\arm{HINGE0}). Part~B ($S = 4$, $P = 4$, $k = 16$, convolution, 28{,}800, 20 seeds): \arm{WIN3\_SLOW16} against \arm{WIN16\_A}, the recurrent gate under Adam with the early hinge window. Parts~C and D ($S = 8$, $P = 4$, $k = 16$, convolution, 43{,}200, 20 seeds): \arm{WIN3\_SLOW\_D8}, \arm{HINGE\_D8} (the recurrent recipe) and \arm{WIN3\_SPLIT\_D8} (the window gate with the split), with a reset-only control (\arm{WIN3\_RESET\_D8}: the Adam-state reset at the same trigger, no copy) on five seeds per machine, printed and not claimed. The perfect gate on two seeds per part was the validity arm, and bound on every one (Part~C at 4800 and 3600--6000). Every CHECK passed on X, including bit-for-bit reproductions of batch~1's window gate, batch~16's \arm{LOCAL3\_SLOW} and batch~17's \arm{W\_SPLIT} (its split at 4800 included); L's file carries the metadata the test writes only after its verification passes. Nothing was cut; X ran its 225 learned-gate records in 30.9\,h of training (its six perfect-gate records came from an earlier start that a container restart ended), L its 231 records in one start (projected 17.3\,h).

\begin{table}[H]
\centering
\small
\setlength{\tabcolsep}{5pt}
\begin{tabular}{@{}llccc@{}}
\toprule
part & arm & X & L & pooled\\
\midrule
A: $S{=}2$, $P{=}4$, $k{=}2$, discovered & \arm{WIN3} & $33/40$ & $32/40$ & $65/80$\\
 & \arm{WIN3\_SLOW} & $33/40$ & $36/40$ & $69/80$\\
 & \arm{HINGE0} (recurrent, Rev.~7's recipe) & $33/40$ & $34/40$ & $67/80$\\
\midrule
B: $S{=}4$, $k{=}16$, bound routed & \arm{WIN3\_SLOW16} & $19/20$ & $19/20$ & $38/40$\\
 & \arm{WIN16\_A} (recurrent, early window) & $15/20$ & $12/20$ & $27/40$\\
\midrule
C, D: $S{=}8$, $k{=}16$, bound routed & \arm{WIN3\_SPLIT\_D8} (the recipe) & $18/20$ & $20/20$ & $\mathbf{38/40}$\\
 & \arm{WIN3\_SLOW\_D8} (no split) & $6/20$ & $9/20$ & $15/40$\\
 & \arm{HINGE\_D8} (recurrent, Rev.~7's recipe) & $0/20$ & $0/20$ & $0/40$\\
 & \arm{WIN3\_RESET\_D8} (control, of 5) & $0/5$ & $3/5$ & $3/10$\\
 & perfect gate (of 2) & $2/2$ & $2/2$ & $4/4$\\
\bottomrule
\end{tabular}
\caption{\texttt{test\_window\_gate}. Pooled claims (primary, exact one-sided McNemar over both machines' pairs): G1, \arm{WIN3\_SLOW16} beats \arm{WIN16\_A}: 13 vs 2 (X 5 vs 1, L 8 vs 1), $p = 0.0037$, SHOWN. G2, \arm{WIN3\_SPLIT\_D8} beats \arm{HINGE\_D8}: 38 vs 0, $p = 3.6\times10^{-12}$, SHOWN. G3, \arm{WIN3\_SLOW} not worse than \arm{HINGE0} by more than two discordant pairs: \arm{HINGE0} only 9, \arm{WIN3\_SLOW} only 11, $d = -2$, HOLDS (a bound, no $p$-value). G4, \arm{WIN3\_SPLIT\_D8} beats \arm{WIN3\_SLOW\_D8}: 23 vs 0 (X 12 vs 0, L 11 vs 0), $p = 1.2\times10^{-7}$, SHOWN. Per machine every claim is SHOWN or HOLDS on both, except G1 on X (5 vs 1, $p = 0.11$; on L 8 vs 1, $p = 0.020$). Bands, pooled: \arm{WIN3\_SPLIT\_D8} RELIABLE, Wilson $[0.835, 0.986]$; \arm{WIN3\_SLOW16} RELIABLE $[0.835, 0.986]$; \arm{WIN3\_SLOW} MAJORITY $[0.770, 0.921]$; \arm{WIN3\_SLOW\_D8} MINORITY $[0.242, 0.530]$; \arm{HINGE\_D8} NEVER $[0, 0.088]$. The pre-registered reading, ``eight streams bind without labels or restarts'' (\arm{WIN3\_SPLIT\_D8} RELIABLE pooled and G2, G4 SHOWN), applies.}
\label{tab:main}
\end{table}

\paragraph{What the runs show (diagnostics, not claims).} The recurrent-gate recipe bound none of 40 eight-stream runs (15 collapsed, 17 non-stream, 8 merged). The window gate alone bound $15/40$, and every one of its 25 failures was a merge of two, three or four streams on one channel. The split raised that to $38/40$: the trigger fired 85 splits in 40 runs (none to three per run), every target labelled on target, and brought the median time to bind forward (X: $21{,}600$ against $27{,}600$; L: $18{,}600$ against $32{,}400$; medians over bound runs). The two failures (X's seeds 555 and 556) ended with two streams sharing a channel. The reset-only control matched the window gate without the split seed for seed on L (3/5 and 3/5, 0 vs 0 discordant) and bound $0/5$ against $1/5$ on X: on these seeds the row copy, not the Adam reset, made the difference (Section~\ref{sec:split}). At four streams, \arm{WIN3\_SLOW16}'s two failures were one unrouted bind (X) and one merge (L); \arm{WIN16\_A}'s thirteen were four position splits, three merges, four unrouted binds and two others. At two streams no arm failed by position; \arm{WIN3}'s failures were key splits (12 of 15), \arm{WIN3\_SLOW}'s key splits and partial streams (6 and 5 of 11), and \arm{HINGE0}'s key splits (10 of 13), with two OTHER and one partial. Slow memory on the window gate at two streams is printed, not claimed: $69/80$ against $65/80$ (7 vs 3, $p = 0.17$).

\paragraph{What is and is not established.} The result is a pre-registered, two-machine, disjoint-seed test of a recipe that uses no stream labels, no restarts and no hinge, and that finds an eight-way partition of the memory, the streams on disjoint channel sets with one channel per stream by largest mass, in 38 of 40 runs where the previous recipe found none. The split needs to know where the key positions are (as the hinge did); the window gate needs the cue within two tokens; and the task is the grouped layout, where it is. Neither limitation is new. The audit's over-prediction for two streams (``$\geq 36/40$ with KEY failures gone'') was not met on X ($33/40$) and KEY failures were reduced, not removed; what was claimed, G3, holds.

\paragraph{Routing by overlap, intervals, and what a replication should expect.} The referee's recount (Section~\ref{sec:referee}) scored the headline runs a second way, by the overlap of the streams' channel sets rather than by each stream's largest channel: in all 38 bound runs the sets are disjoint (largest pairwise overlap $0.047$), and under that criterion G4 is unchanged and G1 strengthens to 16 vs 2. The paired differences carry their uncertainty: G1's four-stream difference is 13 vs 2 ($p = 0.0037$), which is secondary, is not shown on X alone (5 vs 1), and would not survive a family-wise correction over the project's 53 claims (Bonferroni $0.20$), though it survives Holm within the pooled family of six; G3's two-stream difference of two pairs in the window recipe's favour is $+0.025$ with a 95\% interval of $[-0.079, +0.121]$, so ``equal at two streams'' means within eight points either way, not equal. The eight-stream rate, $38/40$, has Wilson interval $[0.835, 0.986]$; a 20-seed replication of a recipe that truly binds $95\%$ of the time reaches $\geq 18$ with probability $0.92$ and $\geq 17$ with $0.98$, and one of a recipe that binds $90\%$ reaches $\geq 18$ with only $0.68$. Section~\ref{sec:replication} reports a replication at $17/20$.

\section{The split: the copy, the reset and the trigger}
\label{sec:split}

The split bundles two actions, a row copy and an optimizer-state reset, and Revision~7's screens had no reset-only control. Revision~8 reported three looks (Table~\ref{tab:copy}, first three rows), each comparing the copy with its reset against the reset alone; they showed that the reset alone is not enough, and the first draft said more than that. Two further looks address it, one of them pre-registered and complete on one machine.

\begin{table}[H]
\centering
\footnotesize
\setlength{\tabcolsep}{3pt}
\begin{tabular}{@{}>{\raggedright\arraybackslash}p{0.13\linewidth}>{\raggedright\arraybackslash}p{0.24\linewidth}>{\centering\arraybackslash}p{0.07\linewidth}>{\centering\arraybackslash}p{0.09\linewidth}>{\centering\arraybackslash}p{0.07\linewidth}>{\centering\arraybackslash}p{0.06\linewidth}>{\raggedright\arraybackslash}p{0.2\linewidth}@{}}
\toprule
where & configuration & copy + reset & copy only & reset only & none & paired\\
\midrule
batch 18, S50 (E) & $S{=}8$, $k{=}16$, 260--269; the reset replayed at the recorded split updates & $9/10$ & --- & $2/10$ & $3/10$ & 8 vs 1, two-sided $p = 0.039$\\
session H, S58b & $S{=}8$, $k{=}16$, 460--469; live trigger & $10/10$ & --- & $2/10$ & --- & 8 vs 0, one-sided $p = 0.0039$\\
\texttt{window\_gate} control & $S{=}8$, $k{=}16$, 540--544, 1540--1544; live trigger & --- & --- & $3/10$ & $4/10$ & 0 vs 1\\
session H, S61 & $S{=}8$, $k{=}16$, 480--489; live trigger & $9/10$ & $9/10$; with noise $9/10$ & $4/10$ & $3/10$ & copy only vs reset only 5 vs 0, $p = 0.031$ (ours, not pre-fixed); vs copy + reset 0 vs 0\\
\texttt{split\_copy}, X & $S{=}8$, $k{=}16$, 560--569; live trigger; pre-registered & $\mathbf{10/10}$ & --- & $\mathbf{3/10}$ & $5/10$ & C1 7 vs 0, $p = 0.0078$ SHOWN; C2 reset vs none 0 vs 2, HOLDS\\
\bottomrule
\end{tabular}
\caption{The split's two actions, bound routed. In every look the reset-only arm fired under the live rule or at the recorded split updates, on labelled targets where labelling was recorded; it bound no more often than the arm with no trigger in \texttt{test\_split\_copy} and the \texttt{window\_gate} control, and more often in S61 ($4/10$ against $3/10$, 2 vs 1). S61's pre-fixed readings ``the copy alone suffices'' and ``the reset adds nothing'' are one-pair bounds on the copy \emph{with} noise against the full split (0 vs 0), and both held; the exact copy (no noise, no reset) enters only the reading ``the noise matters'' ($\geq 4$ vs 0), which did not apply (0 vs 0), and the exact copy against the full split (0 vs 0) was not a reading. \texttt{test\_split\_copy}'s L half has not run; its pooled claims and readings wait for it.}
\label{tab:copy}
\end{table}

\paragraph{The pre-registered look (\texttt{test\_split\_copy}, X).} Specified (\texttt{specs/test\_split\_copy.md}) before it was written, with three claims pooled primary and per machine secondary. On X (seeds 560--569, 54 runs; every CHECK passed except two whose records another CPU wrote, printed rather than asserted, one of them (130) differing; 14.7 h from the test's commit to the result's, the runs about 12 h on four workers): C1, the split beats the reset alone, SHOWN, $10/10$ against $3/10$, 7 vs 0, $p = 0.0078$; C2, the reset alone is not better than no trigger by more than two pairs, HOLDS, $3/10$ against $5/10$, 0 vs 2 (every seed the reset bound also bound without it); C3 is Section~\ref{sec:memory}'s. All 22 splits and all 27 resets were on labelled targets (printed, not claimed); the reset never uses its target, so what the copy adds, on our reading, is the row and not the aim. The validity arms bound 2/2 (the Hebbian perfect gate at 4800 and 3600; the delta perfect gate at 1200 and 1200). The pooled claims and the pre-fixed readings (``the copy, not the reset'' if C1 is SHOWN pooled and C2 HOLDS) wait for L's ten seeds; the readings are pooled only, and on X alone their conditions hold as secondary claims.

\paragraph{The copy alone (session H, S61).} With the exact copy as its own arm, the three forms of the split bound the same nine seeds and missed the same one (483, two streams sharing): the copy with noise and reset $9/10$, the copy with noise $9/10$, the exact copy $9/10$; the reset alone $4/10$, no trigger $3/10$. All 83 firings across the four trigger arms were on labelled targets, a channel holding at least two streams and an empty one; the 57 in the three copying arms copied the row, the 26 resets did not. That two identical rows, their symmetry broken only by their Adam moments and the data, separate merged streams as well as two noisy ones is what 0 vs 0 on ten seeds is consistent with, not a shown equality. The exact copy is therefore the split's form in every later screen and in the next main-line test; the pre-registered test of the copy alone against the full split is to follow.

\paragraph{The constants (session H, S79).} The split's constants had never been varied. With the exact copy as the base ($9/10$, 21 firings, all on target) and one change at a time on ten seeds: a cap of one split, $5/10$ (BASE only 4, CAP1 only 0; ``sensitive''); checks every 1200 from 2400, $7/10$ (2 vs 0; ``inconclusive''); a cap of six, $9/10$; a 16-sequence probe, $10/10$; a rise threshold of $0.05$, $10/10$ (all ``robust''). The pre-fixed reading ``the extra splits are harmless'' did not apply, because one of the cap-of-six arm's seven extra firings was not on a labelled target (the run bound anyway). Descriptively, and not pre-registered: four of CAP1's five failures spent their single split at updates 4800--12{,}000 while the probe read $0.09$--$0.50$ (per-stream accuracy is recorded only at 4800 and 9600; the seed that fired at 12{,}000 had two streams above $0.99$ at 9600), and the fast-check arm's three failures had used all three splits by 7200--10{,}800. The plateau rule fires during the flat start, and a split spent there is one fewer for the merge that comes later; the fast check binds sooner when it binds (median 14{,}400 against 24{,}000) and fails that way when it fails. Our recommendation from this, not a reading: a replication should state the cap and the check schedule (the fast-check arm moved the interval, the first check and the gap together, so the three are not separated), and can state the probe size and the rise threshold loosely. Two trigger variants that follow from this were screened on ten fresh seeds (session~H's S84, seeds 2300--2309; the oracle $2/2$). A guard that forbids splits while the probe is below $0.5$ binds $9/10$ against the base's $9/10$ (0 vs 0); it fires its 17 splits at probe accuracies of $0.50$--$0.91$, where the base fired 7 of its 24 at $0.06$--$0.16$, and it locks out the one seed whose probe never reaches $0.5$, which ends MERGED in every guarded arm. Checks every 1200 with the guard bind $9/10$ (S79's unguarded fast check, on other seeds, bound $7/10$) with a median transition 1200 earlier than the base's, not the 2400 the pre-fixed reading ``faster without the cost'' required. A trigger that reads the 2048-query held-out evaluation instead of the probe binds $10/10$ (1 vs 0; the reading ``the evaluation is the better probe'' applies by its rule, on a single discordant seed, where the probe read $0.000$ then $0.047$ and the evaluation $0.066$ then $0.068$, so, on our reading, the probe's noise decided whether the first check fired); the guard with fast checks and the evaluation together bind $9/10$ (0 vs 0), and no arm read ``costs binding''. These seeds hardly tested the flat start, since, on our reading, the base kept enough splits for its merges after spending some; the guard changes when the splits fire and, on these seeds, not whether the run binds. Session~F's S81 is the same screen on delta channels.

\paragraph{The probe can be wrong about a plateau.} The replication of Section~\ref{sec:replication} lost one of its three failures to the probe: on its seed 914 the 64-query probe read $0.953$--$1.000$ at every check from 16{,}800 on while the 2048-query evaluation sat at $0.871$--$0.887$ with two streams merged, so the rule never fired again. A 64-query probe near $0.9$ has a standard error of about $0.04$. The referee then found that the main line's own fixed probe is unbalanced across streams (5, 5, 6, 8, 10, 12, 5 and 13 queries for streams 0--7) and identical for every seed, so that a merge of two under-sampled streams costs the probe about $0.08$ where the balanced evaluation loses $0.125$: the direction that blocks a split, and a systematic bias rather than noise. The main line's records store the probe's reading at every check beside the evaluation's curve, and the referee read them against each other: 6 of the 353 checks in the split arm's 40 runs read the probe at or above $0.95$ with the evaluation at $0.865$--$0.88$, in 2 of those the probe alone blocked a split, and both fired at the next check and bound routed; the reset arm had none. The next main-line test prints every check's probe reading beside the held-out accuracy, and the trigger arms above are the remedy under test.

\paragraph{What does not replace it (session H).} Two devices the literature uses against router collapse were run beside the split and the reset at four streams, $k = 16$ (40 paired seeds), where every arm but one sat at the ceiling (NONE $38/40$); two of its three readings did not apply and the third, ``the split is the reset'', held by its rule (0 vs 0) and was printed uninformative by its own qualifier. A rule fixed beforehand then carried the two best-scoring devices, the reset and write-gate noise (it ranked the four devices only; the no-device arm's $38/40$ outscored write-gate noise's $37/40$), to eight streams (10 seeds). Gumbel$(0,1)$ noise on the gate's logits during training, Raven's device for a router whose reads are dense \cite{raven}: on both the read and write gates it collapsed routing even at four streams (33 of 40 runs merged); on the write gate alone it bound $37/40$ at four streams with a later transition (median 6000 against 3600) and a lower margin, and at eight streams it bound $0/10$, with up to six streams sharing a channel (10 vs 0 against the split; the screen had no no-device arm there). A Switch-style balance loss, as Mixture-of-Memories uses \cite{mom, switch}: $36/40$ at four streams, routed by stream, but in 34 of those 36 runs the streams were mapped to different channels at key positions than at value positions and the routing margin fell from $0.87$ to $0.46$; on our reading it changed the routing's shape, not its count, and it was not carried to eight streams. In the other arms with a full set of runs the query's read gate follows the write map (session H's S62, descriptive: in 32--37 of the 38--40 bound runs per arm at four streams, and $12/12$ at eight). At two streams ($k = 2$, no convolution, ten seeds), a learned gate temperature with a floor ran to the floor in every run, and a stable-max softmax read ``no change at $n = 10$'' by its rule ($7/10$ against $8/10$). Slope annealing of the gate (S63) and a hard, one-hot write gate from update 2400 both read ``does not help'' ($9/10$ against $10/10$; $0/10$, 0 vs 10, KEY 3, MERGED 4, OTHER 3). With the write one-hot the gate gets gradient only through the read, by construction; that the writes' stream labelling therefore never forms is the session's description, and four of the ten runs did form a partial one (MERGED needs a margin of $0.25$).

\paragraph{The split needs an idle channel.} With no spare channel (batch~18, S51, $S = k = 4$), the trigger once fired on a correctly routed map, found no empty channel, overwrote a live one and merged the run. Session~H's S78 then ran the exact copy at $k = S = 2$ (Section~\ref{sec:keysplits}): the copy resets the gate to uniform and, on the eight firings of the seeds that ended in a key split, the same split is back by the next check (two of the eight are read from the end of the run); one seed had no key split when the copy came and discovered. That the copy works at eight streams because the channel it copies onto is empty, and has nothing to offer where every channel is busy, is our reading of these screens. Spare channels remain the default, and, our recommendation, the trigger should not fire when its target already carries key mass.

\section{Muon is not an ingredient}
\label{sec:muon}

Revision~7 reported, from one machine, that under Muon the recipe ``binds four streams sooner''. The Muon-recipe test (\texttt{test\_muon\_recipe}, the first on disjoint seeds) asked three questions. M1, the early hinge window under Muon at two streams: \arm{WIN\_M} $79/80$ against \arm{SLOW\_M} $69/80$ pooled, 10 vs 0, $p = 0.001$, SHOWN (on L 6 vs 0, SHOWN; on X 4 vs 0, $p = 0.0625$). M2, Muon binds four streams by update 4800 more often than Adam: 12 vs 8 pooled, $p = 0.25$, NOT SHOWN; per machine it was shown on X (7 vs 1, $p = 0.035$) and reversed on L (5 vs 7). Revision~7's result had come from L (7 vs 0 on its seeds; X's cut arm gave 5 vs 1), so on fresh seeds it was confirmed on the other machine and reversed on its own; the machines disagree, by seed block or by CPU (the design cannot separate them), and pooled it is not shown. M3, the same at eight keys: 7 vs 5 over 30 pairs (X's time rule cut its Adam arm to ten seeds), $p = 0.39$, NOT SHOWN. The test's reading, ``the Muon recipe is reliable at the working configurations'', was not met: \arm{WIN16\_M} was MAJORITY pooled ($33/40$). Batch~18 then ran the window recipe with the split under Muon at eight streams: $9/10$, the same as under Adam (1 vs 1), and $9/10$ against a fresh Muon reference without the split at $2/10$ (7 vs 0, two-sided $p = 0.016$). The eight-stream result is the recipe's, not the optimizer's. Muon's bf16 Newton--Schulz step does not reproduce across processors, which cost the project paired comparisons; nothing in the main line now needs it.

\section{The recurrent gate's eight-stream failure, closed}
\label{sec:recurrent}

Revision~7's reading was that the recurrent gate's state stops carrying the stream when its recurrence becomes expansive, and it left a cap on that recurrence being screened. Batch~15 (S41) ran the cap at $0.5$ on the spectral norm of $W_h$ under Muon, alone, with a previous-token input, and with the previous-token input and no recurrence, and the capped previous-token input under Adam: the three Muon arms bound $0/10$ and the Adam arm $0/5$. With the look-back the capped gate's state kept the stream at key positions throughout (decodability $0.93$--$1.00$) and still ended with the eight streams on three to seven channels; final accuracy rose in every pair (median $0.61$ against $0.36$) but no run bound. Batch~16 added a frozen reservoir recurrence ($0/10$), decay $0.98$, a tenth of the rate on $W_h$, and a higher learning rate: none bound (the perfect gate binds at decay $0.98$, so the decay is valid). The window gate then bound where all of these had not. Two corrections to Revision~7 follow: the gate's state must carry the stream at the \emph{write} position, where CTX is two tokens back, not only at the read, where it is one back, so the sentence ``the gate needs only one token of memory'' was wrong (the audit noted it; batch~16's S47 measured the write positions, which S38--S41 had not); and S41's no-recurrence arm, whose input is $(v_t, v_{t-1})$, cannot by construction see the stream token at value positions, so its ten key splits were a property of that window and not evidence about recurrence (audit, Section~2.3). The cap, as specified, also acted from update 1 ($\sigma_{\max}(W_h) \approx 1.08$ at initialization), not only after the crossing. The diagnosis stands as a description of the recurrent gate; the remedy was to remove the recurrence.

\section{Two-stream key splits}
\label{sec:keysplits}

At two streams with two channels and no convolution, the window gate fails by routing on key identity: its key splits form within the first 600 updates, while the memory is still at $10^{-4}$ (batch~17, S49). Slow memory reduces them (\arm{WIN3} KEY 12 of 15 failures pooled; \arm{WIN3\_SLOW} KEY 6 of 11), and batch~18 screened two further devices on seeds 160--199: two spare channels ($k = 4$: $36/40$ bound routed, with four key splits still formed by update 325) and the hinge's key term on updates 1--2400 ($36/40$ discovered, one KEY failure, but three ``STREAM-PARTIAL'' failures with $\eta^2$ by key within $10^{-3}$ of $0.5$, the signature of a hard two-of-four key pattern; the key-first classifier relabels the one at $0.500004$ KEY and leaves the two just below). The configuration is the project's hardest relative to its size and the one with the least headroom ($65$--$69/80$ for every arm); it is also the only one where the window recipe's margin over the hinge recipe is small (two pairs pooled, $[-0.079, +0.121]$) and where the hinge recipe won a screen (batch~16: $34/40$ against $31/40$ on seeds 160--199).

\paragraph{The split does not repair a key split (session H, S78).} Session~H ran the exact copy at $k = S = 2$ on 20 seeds (1300--1319), with the main line's noisy form and the two-spare-channel recipe beside it: the baseline discovered $15/20$, the exact copy $16/20$, the copy with noise and reset $17/20$, four channels $15/20$; no arm beats the baseline (best 2 vs 0, $p = 0.25$), and the pre-fixed reading is ``neither device''. The firing records say why. After $W_g[c_0] := W_g[c^*]$ the two rows are equal and the gate is uniform, every $\eta^2$ at $0.00$; the same key split had re-formed exactly (KEY/stream $\eta^2$ $1.00/0.00$) after all eight firings on the three seeds that ended KEY under the classifier of record (six read at the next check, 2400 updates later; two, fired at the last check, read from the end of the run), and the half-key pattern after all three on a fourth. The one rescue was a seed whose gate was not yet a key split when the copy came. With noise and the Adam reset the key split re-formed only partly on two seeds and fully on one. Our reading: at eight streams the copy creates a near-tie between a crowded channel and an empty one; at $k = S$ it only resets the gate to where the gradient takes it again. Session~H's reading of the mechanism, that $W_g$'s untouched Adam moments still point along the old split, was not measured and is not a reading.

\paragraph{Prevention, not repair (session H, S85).} Two devices against the key split were screened on 20 fresh seeds (2310--2329; the oracle $2/2$ at 1200), under the key-first classifier: a gate held at $10^{-4}$ for the first 1200 updates, and the hinge's key term on updates 1--2400 (batch~18's device, Section~\ref{sec:keysplits} above). The key term discovers $\mathbf{20/20}$ against the baseline's $15/20$ (5 vs 0, one-sided $p = 0.031$) with no key failure (the baseline has four), and its pre-fixed reading ``prevention works'' applies; the slow gate discovers $14/20$ (0 vs 1) with four key failures of its own, three on the baseline's seeds, reading ``prevention not shown; the key split only moves''; the two together discover $19/20$ (5 vs 1, $p = 0.109$), ``not shown'', the one failure (seed 2317) a run that each device alone discovered: under both together the gate was still unrouted at 2400 and the split formed after the window ($\eta^2$ by key $0.30$ at 3600, $0.77$ at 4800) without recovering, which session~H reads as the slow gate pushing the gate's learning past the hinge's window, though its hold ends at 1200 and the record shows only the combination failing. Descriptively, from the gate's $\eta^2$ by key every 25 updates to 1200 and every 100 to 2400: the baseline's onsets after update 0 lie at updates 25--125 (seven seeds) and the slow gate's at 275--575 (three seeds, two of them seeds the baseline also split and one it discovered), so, on our reading, the gate's learning rate is not what lets them form; under the key term $\eta^2$ by key stays below the term's threshold of $0.2$ at every measurement from update 25 to 2400 (maximum $0.195$ over the 20 runs; the term fired on 1--11 batches per run), a split can start after the window ($0.21$--$0.88$ at 3600 on seven seeds) and every such run recovers, the stream routing having formed first; the price is a later transition (median 4800 against 2400). At initialisation the near-uniform gate's small variations are organised by key on five of the twenty seeds ($\eta^2$ up to $0.84$ at update 0), the audit's near-uniform case, which the baseline's discovering runs shed by update 600. This is the result of session~G's S82 at the other place the key split forms (Section~\ref{sec:far}): the key term, which reads key positions and key identities during the slow phase, prevents the split that the copy did not repair (S78). Session~H proposes the pre-registered test, the window recipe with and without the key term at two streams on disjoint seeds under the key-first classifier. Four of S78's seventeen unbound runs are STREAM-PARTIAL under the classifier of record and KEY under \texttt{fail\_class\_v2}, and one of S85's; no reading changes.

\section{Scale}
\label{sec:scale}

Batch~19 (E, running) asked whether the recipe's constants are tuned to $N = 256$, $D = 32$, three layers. At four streams with sixteen channels (seeds 270--289, the oracle binding at every size and depth), with every constant unchanged: $N = 512$ and $1024$ at $D = 32$, and $D = 64$ and $128$ at $N = 256$, bound routed $10/10$ each against the reference's $9/10$, all with median transition 3600; 2, 4 and 6 applications of the shared layer bound $10/10$, $10/10$ and $9/10$ against $10/10$ for three, with six layers slower (median 4800). Both readings were ``scale-free'' by their pre-fixed rule (every $d \geq -1$), at four streams only; with ten seeds the screen rules out a large drop, not a small one. Eight streams and larger key loads at the larger widths are running.

\section{The memory rule}
\label{sec:memory}

\paragraph{What the literature holds.} Twelve papers on routing, memory and context were read in full (\texttt{docs/reading/}): the test-time-regression framework \cite{ttr}, Gated DeltaNet \cite{gdn}, the Tiny Recursive Model and recurrent-depth language models \cite{trm, recdepth}, Flesch et al.'s Hebbian context gating \cite{flesch}, Active Dendrites \cite{dendrites}, Context-Gated Associative Retrieval \cite{cgar}, Mixture-of-Memories \cite{mom}, Raven \cite{raven}, a study of content-based routing \cite{routing}, the Hierarchical Multiscale RNN \cite{hmrnn} and Nested Learning \cite{nested}. None of the twelve carries an inferred context cue across distance without labels under a sparse end-of-sequence loss (the Hierarchical Multiscale RNN infers its boundaries from a dense likelihood, as our note records): each hands the context over (a task identifier, a prototype readable off every input, a clamped context layer), carries it with softmax attention, or routes by token type. The study of content-based routing is the sharpest statement of why our header-layout runs failed: a router over a frozen state-space model's summaries reaches about $29\%$ of the right positions even with the route supervised, over Fourier-mixed representations $0.9\%$, and $98$--$99.7\%$ once an upstream, jointly trained attention or mixing layer has written the context into each token's representation. Flesch et al. fit a leaky integrator of the context cue to human learners and show that slow integration of interleaved contexts merges them, and that inside a block the loss exerts no pressure to use a constant cue. Our window gate reads three raw embeddings; our fading-memory gates pooled everything; our fixed header layout is positionally solvable. Those are the failures the papers' analyses, read as analogies, predict.

\paragraph{BDH-CQ.} Pathway's follow-up \cite{bdhcq} describes its memory as ``linear correction rules on $S$'' with linear attention as the special case $S_t = S_{t-1} + U_\theta(D_t)$, and its reasoning as $H_{r+1} = F_\theta(H_r, S_K)$ over a frozen memory; ``dimensions, exact update rules, and implementation details remain proprietary''. Our reading (\texttt{docs/reading/}, an inference from the sentences quoted above, not a finding, and one reading among several, since the quoted sentence counts plain linear attention, which does not read $S$, among such rules): in the taxonomy of \cite{ttr} and \cite{nested}, the rules that are affine in $S$ and read $S$ before writing are the delta-rule family (DeltaNet, Longhorn, Gated DeltaNet, DeltaProduct), and the loop matches the recursive-refinement models \cite{trm, recdepth}; the recurrent-depth model re-injects its input embedding at every step to keep the recurrence stable, and the frozen memory $S_K$ in BDH-CQ's loop would be the analogue, which is our analogy and not theirs. Neither is a distant-cue mechanism. The $97.4\%$ Sudoku figure sometimes attached to BDH is from Pathway's blog, on internal data with an unstated protocol, and is not comparable to the Sudoku-Extreme numbers of \cite{trm}.

\paragraph{The premise: one channel cannot.} One Hebbian channel binds the two-stream grouped layout $0/60$ without the convolution and a minority with it at a higher learning rate ($18/40$, $10/20$), and four streams $0/20$ (Phase~IV). Session~F's delta-rule channel (Section~\ref{sec:task}) was the stronger single memory the project lacked: at a fixed decay it binds the two-stream \emph{header} layout $8/10$ (S53, $P = 4$) and the grouped two-stream layout $4/10$. The premise test then ran it where the partition's claim lives (S68, seeds 300--309, 43{,}200 updates): one delta channel binds the grouped layout $2/10$ at four streams and $0/10$ at eight, while both delta perfect gates bind at 1200 and the gated Hebbian recipe binds $19/20$ and $9/10$ on other seeds; the pre-fixed readings ``context-tagged keys scale'' ($\geq 7/10$ at eight) and ``the partition's advantage is in the number of streams'' ($\geq 7/10$ at four and $\leq 3/10$ at eight) both read ``neither''. The main line's pre-registered version, \texttt{test\_split\_copy}'s C3, is SHOWN on X: the recipe $10/10$ against one delta channel $0/10$ (10 vs 0, $p = 0.00098$; final accuracy $0.08$--$0.76$), one Hebbian channel $0/10$ (printed, not claimed) and the delta perfect gate $2/2$ at 1200; the pooled claim and the reading ``the partition is necessary at eight streams against a single delta channel'', which is what would make the channel's failure the missing partition rather than the delta rule, wait for L.

\paragraph{The forget gate needs the delta rule.} On the fixed eight-key header layout a single delta channel with a learned forget gate (Raven's routed decay, or Gated DeltaNet's) bound $9/10$ and $8/10$ against $3/10$ at the fixed decay (S59), and, post hoc, had learned to partly clear its memory at the context token (median $\alpha$ of $0.18$--$0.26$ at CTX in the first layer, $0.72$--$0.74$ in the third, $1.00$ at keys and values). The same decays on a single \emph{Hebbian} channel, same seeds and batches, bound $0/10$, $0/10$ and $0/10$ (S69): the Hebbian channel does not learn the clearing (median $\alpha$ at CTX $0.975$ and $0.988$ in the first layer), and its final accuracy was lower with the learned decay ($0.17$ against $0.22$; a description of ten runs).

\paragraph{The delta recipe keeps eight streams (session F, S76).} Revision~8 reported that the window-gate recipe on delta channels stalled at eight streams ($3/10$ against $9/10$, S54(c)) with a perfect gate that itself took 22{,}800--28{,}800 updates, and concluded that the Hebbian write stays on the main line. S76 repeated the comparison with the write strength fixed at $\beta = 1$ and the exact-copy split. Its first phase ran the perfect gate on four seeds: with L2-normalised keys the delta oracle binds all four at the first evaluation (1200) at $\beta = 1$ and at 2400 at $\beta = 0.25$, at $\beta = 1$ sooner than the Hebbian oracle on every seed (2400--6000; at $\beta = 0.25$ it ties the Hebbian oracle on two of the four); with raw keys at $\beta = 1$ one seed diverged and one bound late (16{,}800), and at $\beta = 0.25$ all four bound, one of them at 7200. The pre-fixed readings for the slowness (the overwrite; the key normalisation) both read ``neither'', because no L2 arm was slow. Its second phase ran the recipe, paired, on ten seeds: on delta channels $9/10$ bound routed against $10/10$ on Hebbian channels (1 vs 0), sooner on eight of the nine seeds both bind (median 14{,}400 against 24{,}000); the reading ``the delta recipe keeps eight streams'' applies. Six of the twelve splits on delta channels were not on labelled targets: at each the map had already merged; in five of the six the split copied a channel that held no stream or one, and at the sixth (seed 1314, update 9600) it copied a merged channel onto an occupied one; mostly single splits in runs that then bound. That the plateau rule's target selection is tuned to the Hebbian trajectory is our reading of this. S54(c) differed from S76 in the learned $\beta$, the noisy split with a reset, and the seeds; the oracle comparison holds the routing fixed, which makes the learned $\beta$ the likely cause of the earlier stall, consistent with its collapse toward zero in S53 and S54(d), but the screen did not separate the two and that attribution is post hoc.

\paragraph{\ldots and two and four streams (S80).} The same recipe on delta channels, paired with Hebbian on ten seeds each: at two streams ($k = 2$, no convolution) $10/10$ against $10/10$, delta at the first evaluation on every seed against 2400--6000; at four streams ($k = 16$, with the convolution) $10/10$ against $10/10$, both at 3600 on every seed, no split firing in either arm. Both readings are ``keeps it''. Both arms sit at the ceiling, so ten seeds can only detect a large loss: at the main line's rates ($0.86$ and $0.95$) a Hebbian arm loses about $1.4$ and $0.5$ runs in ten, and ``keeps it'' means that no seed failed on delta, not equivalence at the main line's power.

\paragraph{What the delta stack carries.} Session~F's one-seed observation in Revision~8, keys at the third layer tagged by stream, did not generalise. Reruns of 30 recorded single-channel runs with a linear probe at the end (S77; every curve bit for bit) found the stream linearly readable from the residual at layer~3 in one of sixteen bound delta runs on the fixed header layouts (medians at key positions $0.57$--$0.67$), and the two S53 runs that did carry it at value positions ($0.92$--$1.00$) were the two that failed to bind; the Hebbian channel's medians are $0.50$--$0.53$ at key positions (single runs up to $0.60$). Our reading: keys can be near-orthogonal in the 256-dimensional code without a stream direction in the residual. On the grouped two-stream layout (S80) the delta stack does carry the stream at \emph{key} positions, from layer~2 on ($0.98$--$1.00$ at layer~3, median $1.00$; Hebbian $0.57$), and not at value positions ($0.54$--$0.65$), in the perfect-gate arm as well; session~F reads this, post hoc, as adjacency to the context token that precedes each key rather than a tag carried to the write. On the randomised header layout (Section~\ref{sec:far}) the in-run probe finds the stream at key positions in most of the delta solutions that bind without routing and in no routed one.

\paragraph{What this changes.} ``The Hebbian write stays on the main line'' is withdrawn. With the write strength fixed, the delta write is not shown worse than the Hebbian write at two, four and eight streams on screens of ten paired seeds (``keeps it'': $10/10$, $10/10$ and $9/10$ against $10/10$ each, one-sided bounds at the ceiling and not parity), binds sooner at two and eight, binds the randomised header layout $14/20$ against $1/20$ for the window baseline (Section~\ref{sec:far}), tolerates an unrouted first layer (S71), and as a single channel binds $2/10$ and $0/10$ at four and eight streams (S68, whose reading of that name did not apply; the pre-registered support for the partition is C3 on X). The question is a main-line one again, and \texttt{test\_delta\_recipe} puts it (specified in \texttt{specs/}, reviewed by the referee before any run, amended once on its review, and written, with the test file checked against the amended spec; not run): the exact-copy recipe and the recipe without a trigger on delta and on Hebbian channels (four arms), 20 seeds per machine. D1 is the pooled discordant difference between the two split arms, labelled ``delta binds more often'' ($d \leq -3$), ``delta not shown worse by more than 2 pairs'' or ``delta not shown non-inferior'' ($d \geq 3$), printed beside a paired rate difference with a 95\% interval and the bound's operating characteristics (a ten-point deficit passes it 22--29\% of the time at 40 pairs); D2 and D4 test the split against no trigger on each write (exact one-sided McNemar) and read ``helps'', ``not shown'' or, when the no-trigger arm binds $\geq 36/40$, ``no room''; D3, secondary, is that delta binds sooner. The referee's review, which found the first draft's readings ``interchangeable'' and ``Hebbian-specific'' each readable two ways, is in \texttt{docs/review\_2026-10/}. Raw keys are excluded by their numerical fragility ($\beta\|x\|^2 > 2$ makes $I - \beta x x^\top$ expanding, which is session~F's post hoc reading of the one divergence in eight raw-key oracle runs), and the trigger on delta channels is the subject of F's S81.

\section{Distant cues}
\label{sec:far}

The window gate sees three tokens, so on the header layout it cannot see the cue from the third pair of a block on. Revision~8 left the distant-cue problem open with three proposals from the audit: probe where the cue is in the residual stream, add a register that overwrites on the cue and holds otherwise, and let a per-layer gate read the residual stream. Session~G ran all three, in two rounds, and the second round changed the answer.

\paragraph{The cue is not in the stack (S55).} A logistic-regression probe from the residual stream entering each layer to the stream, at key and value positions, on the trained window-gate stacks of the eight-key header layout: $0.55$--$0.58$ at layers 2 and 3 and at the output, from about $0.70$ at initialisation (where, on our reading, the decayed sum of the block still carries the header). Training removes the cue (our reading: the gate reads only the window, so nothing asks the stack to keep it), and the perfect gate's stack carries nothing either ($0.53$--$0.56$ over the three sites). The pre-fixed reading ``absent'' applied; that a gate over the residual stream therefore had nothing to read, as the routing study would predict \cite{routing}, is our reading of it.

\paragraph{The register at layers 2--3 fails on Hebbian memory, and the reason was layer 1 (S56, S57, S71).} With the register gate at layers 2 and 3 and the window gate at layer 1, on the fixed eight-key header layout, every Hebbian arm bound $0/10$: the baseline, the residual gate alone, the register with and without a stop-gradient, and the register with a 5\% labelled nudge. The nudged arm is the informative failure: its register writes a clean latch in eight of the ten runs (medians: $\beta^R$ $0.96$ at the context token against $0.02$ and $0.05$ at keys and values), its gates route one-to-one with margin $1.00$ at layers 2 and 3 in the same eight, and its accuracy plateaus at $0.20$--$0.44$ (the lower of the two runs without a clean latch) with the two streams within $0.03$ of each other in most runs, no better than the unrouted baseline ($0.28$--$0.42$); the reading ``architecture insufficient'' applied. A perfect gate at layers 2--3 alone bound $0/3$ where the full perfect gate bound one of two seeds, one paired seed of evidence. On delta channels the same architecture binds: the nudged register $8/10$, routed $7/10$, transitions 3600--6000; the unnudged register $6/10$ bound against the baseline's $2/10$, which met the screen's own rule for ``discovered'' on delta (6 against 2 with no unbound key failure; 4 vs 0 by McNemar, which was not its rule), but only two of its six routed, three binding through key splits. The Hebbian perfect gate itself is marginal on this layout at 24{,}000 updates (3 of 6 seeds over the three screens), so those zeros are weak. Session~G then put the layer question on ten paired seeds of the randomised header layout, whose perfect gates bind at 1200 (S71): on Hebbian memory the perfect gate at all three layers binds $10/10$ (median 1200, range 1200--4800), at layers 2--3 alone $0/10$, the runs' accuracies $0.31$--$0.54$ in two clusters near $0.32$ and $0.53$ and the per-cell medians over runs $0.46$--$0.51$ in every (stream, block, pair) cell, with layers 2--3 routed perfectly, and at layer~1 alone $9/10$, slowly (median 16{,}800); on delta memory the perfect gate at layers 2--3 alone binds $10/10$. The pre-fixed readings ``layer~1 is the block'' (10 vs 0, $p = 0.00098$), ``layer~1 alone suffices'' and ``delta tolerates an unrouted layer~1'' all apply. The interpretation, which is ours and not a reading, is that an unrouted Hebbian memory at the first layer mixes the streams before the routed layers read it; the first design kept layer~1 on the window gate, and that choice is what sank its Hebbian register; the interpretation crosses two tasks (S56's fixed eight-key layout; S71's randomised one) and two gate inputs (S56's register gate read the raw embedding, REG3's reads the window).

\paragraph{The register at all three layers carries the cue; discovery is the gap (S72, S73).} One register gate shared by all three layers (REG3, Section~\ref{sec:task}), on the randomised header layout, 20 seeds, readings on ROUTED-BOUND. With the 5\% labelled nudge it carries the cue every time: $\mathbf{20/20}$ on Hebbian memory (transitions 2400--8400) and $\mathbf{10/10}$ on delta (all at 1200), one-to-one with margin $1.00$ at every layer. Without the nudge it finds the solution on $4/20$ seeds on each memory, against $0/20$ for either baseline, each discovering run routed perfectly; both screens are one discordant pair short of their $p < 0.05$ rule (4 vs 0, $p = 0.0625$), so S72's pre-fixed reading is ``not discovered'' and S73's ``discovered (register)'' does not apply, and 8 vs 0 over the 40 pairs is a description, not a test, since the two screens share seeds but not models. Two other readings of S73 did not apply either, ``the delta stack carries the cue'' (the baseline $7/20$) and ``the gate is not needed'' (one channel $7/10$). The S73 reading was fixed on ROUTED-BOUND; on the outcome named DISCOVERED the same arm is 5 vs 0 ($p = 0.031$), one run binding and separating its values with a margin of $0.66$, which is why Section~\ref{sec:task} fixes the names. Descriptively: in every routed unnudged run $\beta^R$ at the context token rose above its value at keys and values (Hebbian $0.32$--$0.64$ against $0.16$--$0.24$; delta $0.36$--$0.58$ against about $0.2$), and in the runs that failed the gate had locked into a key split (14 of the 16 runs that did not route on each memory; on delta two of the 14 bound) with $\beta^R$ at the context token not separated above keys and values in 27 of the 28 such runs. The unbound Hebbian runs sit near chance between the streams (per-cell medians $0.46$--$0.55$; flat across cells by the screen's rule in about half of them). The other devices of the screen: the residual gate alone routes $0/10$ on delta ($5/10$ bound, three by key splits); the HM-RNN latch with slope annealing holds an exact latch when nudged ($9/10$ routed, $z = 1$ at the context token and $0$ elsewhere) and, unnudged, goes to $z = 1$ at every role by update 1200 in eight of ten runs and key-splits in nine of its nine unbound runs (on the fixed layout, S57, it had gone to $z = 1$ in every run and key-split $3/10$); one delta channel binds the layout $7/10$ with no gate. The in-run probe finds the stream in the residual of no routed run; among the 31 non-routing delta solutions it finds it at key positions in 18 and at both key and value positions in 13 (the single channel's bound runs at $0.93$--$1.00$ at keys; the half-routed baseline's at $0.86$--$1.00$, also at keys): on delta memory the stack can learn stream-tagged keys, and when it does the gate does not need to route, which is a description of a $2\times2$ the screen fixed as descriptive.

\paragraph{Holding the key split off lets the write appear (S82).} Most of the runs that fail to discover end in a key split (28 of 32 in S72 and S73), formed in the records before any write at the cue (Section~\ref{sec:keysplits}), so session~G screened three devices against the key split on 20 fresh seeds (2100--2119), Hebbian memory, with the baseline, REG3 and a nudged pair for validity ($2/2$, with the perfect gate $2/2$). The hinge's key term of batch~18, a penalty on the shared gate's dependence on key identity at key positions during updates 1--2400 (REG3\_KEYHINGE), binds routed $\mathbf{20/20}$ against the baseline's $0/20$ and REG3's $2/20$ on these seeds (20 vs 0 and 18 vs 0; the readings ``discovered'' and ``helps'' both apply, and the arm has no key failures and $\beta^R$ at the cue above keys and values in every run); on delta memory, descriptive by the screen's docstring, the same arm binds routed $10/10$ against $0/10$ and $2/10$ (bound: $2/10$ and $4/10$). A gate held slow for the first 1200 updates binds $0/20$ (0 vs 0 against the baseline; ``hurts'' did not apply, 0 vs 2 against the plain register): the records show the key split delayed ($\eta^2$ by key above $0.5$ at 1200 in 2 of 20 runs against 16 of 20) but formed (KEY 18 of 20) and the write at the cue separating in only 1 of 20, which we read as the slow gate delaying the write along with the split; a gate without the raw window input binds $5/20$, ``discovered'' exactly at its threshold (5 vs 0), with a register that writes at nearly every token. What the firing records show, as a description: the near-uniform gate's small variations are organised by key from initialisation ($\eta^2$ by key $0.89$ at update 1), the hinge holds that at $0.05$ from update 300, the register's write at the cue separates from its write at keys and values during the hinged phase ($0.12$ at every role at 1200; $0.48$ against $0.20$ at the end), 18 of the 20 bind at 3600, the first evaluation after the hinge stops, and the key split does not return in the 2400 updates each run is observed after its transition. The hinge reads the task's key positions and key identities, as the split reads key positions: it is a training-time prior over the first 2400 updates and not an architecture, and in a language model key identity would be token identity, which a gate must in general be allowed to use.

\paragraph{What this changes.} On a screen of twenty seeds, the distant-cue task has a solution without stream labels: a register that overwrites at the cue, a gate over the register, the window and the residual at all three layers, and a hinge that keeps the gate off key identity while the memory is slow. Its capacity was never the limit (the nudged register, $20/20$); its discovery is $20/20$ with the hinge and $2$--$4/20$ without. Session~G proposed the pre-registered test and is drafting its specification (S83, started before this draft, screens the control and what the hinge must read); it is the register test proposed to follow \texttt{test\_delta\_recipe}: the hinged register against the baseline, against the register without the hinge, and against the hinge on the window gate without the register (a control the screen did not run), 20 seeds per machine, with a rate claim of $\geq 18/20$ on each. The audit's per-layer gate over the residual stream, run alone, bound $0/10$ on the fixed Hebbian layout (S56, every run a key split, with nothing in the stack to read) and routed $0/10$ on delta (S73; $5/10$ bound, eight key splits); no reading of either screen applied to it, and it is not pursued.

\section{Where the mechanism stands}
\label{sec:synthesis}

\begin{table}[H]
\centering
\footnotesize
\setlength{\tabcolsep}{2.5pt}
\begin{tabular}{@{}>{\raggedright\arraybackslash}p{0.17\linewidth}>{\centering\arraybackslash}p{0.07\linewidth}>{\centering\arraybackslash}p{0.11\linewidth}>{\centering\arraybackslash}p{0.1\linewidth}>{\centering\arraybackslash}p{0.13\linewidth}>{\centering\arraybackslash}p{0.11\linewidth}>{\centering\arraybackslash}p{0.08\linewidth}>{\centering\arraybackslash}p{0.09\linewidth}@{}}
\toprule
setting & perfect gate & recurrent gate, Rev.~7 recipe & window gate + SLOW & window gate + SLOW + split & the same on delta channels$^{d}$ & one Hebbian channel & one delta channel\\
\midrule
$S{=}2$, $P{=}4$, no conv. & $4/4$ & $67/80$ & $69/80$ & --- & $10/10$ vs $10/10$ & $0/60^{c}$ & $4/10^{f}$\\
$S{=}4$, $P{=}4$, $k{=}16$, conv. & $4/4$ & $27/40^{w}$ & $38/40$ & $39/40^{e}$ & $10/10$ vs $10/10$ & $0/20$ & $2/10^{f}$\\
$S{=}8$, $P{=}4$, $k{=}16$, conv. & $4/4$ & $0/40$ & $15/40$; X $5/10^{s}$ & $\mathbf{38/40}$; X $10/10^{s}$ & $9/10$ vs $10/10$ & $0/10$ (X) & $0/10$ (X), $0/10^{f}$\\
\bottomrule
\end{tabular}
\caption{The grouped layout: the main configurations of \texttt{test\_window\_gate} (pooled over X and L, disjoint seeds; bound routed, or discovered at $S = 2$), with \texttt{test\_split\_copy}'s X half and the screens. $^{w}$With the early hinge window (\arm{WIN16\_A}); Revision~7's full-hinge arm bound $35/40$ on shared seeds, routed $15/20$ on X and $16/20$ on L. $^{e}$Session~H's S58, exploratory, 40 seeds. $^{c}$One Hebbian channel without the convolution, every such run on the grouped layout in \texttt{results/} (the blocked and shuffled layouts add $0/30$); with the convolution at learning rate $4\times10^{-3}$ it binds $18/40$ and $10/20$ at two streams (Phase~IV's \texttt{conv\_lr}; Revision~6 reported $34/80$) and $0/20$ at four (\texttt{test\_scale\_axes}, B4s). $^{d}$Session~F's S80 and S76, exploratory, ten paired seeds each, the delta recipe against the Hebbian one with $\beta = 1$, L2 keys and the exact-copy split at four and eight streams (none at two, where $k = S$). $^{f}$Session~F's S53 and S68, exploratory, $\beta = 1$ at the fixed decay. $^{s}$\texttt{test\_split\_copy} on X.}
\label{tab:synthesis}
\end{table}

\begin{table}[H]
\centering
\footnotesize
\setlength{\tabcolsep}{4pt}
\begin{tabular}{@{}>{\raggedright\arraybackslash}p{0.44\linewidth}cc@{}}
\toprule
randomised header layout ($S{=}2$, $P{=}4$, $k{=}2$, 24{,}000; session G, seeds 700--719) & Hebbian memory & delta memory\\
\midrule
perfect gate, all three layers & $10/10$ & $10/10$\\
perfect gate, layers 2--3 only (window gate at layer 1) & $0/10$ & $10/10$\\
perfect gate, layer 1 only & $9/10$ & ---\\
window gate + SLOW (ROUTED-BOUND; bound) & $0/20$ ($1/20$) & $0/20$ ($14/20$)\\
register gate at all three layers, labelled nudge & $\mathbf{20/20}$ & $\mathbf{10/10}$\\
register gate at all three layers, label-free & $4/20$ & $4/20$\\
the same with the key hinge for updates 1--2400 (S82, seeds 2100--2119; baseline there $0/20$, register $2/20$) & $\mathbf{20/20}$ & $\mathbf{10/10}$\\
the same with a slow gate for 1200 updates / without the window input (S82) & $0/20$ / $5/20$ & ---\\
residual gate at all three layers, no register & --- & $0/10$\\
HM-RNN latch, label-free / nudged & --- & $0/10$ / $9/10$\\
one channel, no gate (bound) & --- & $7/10$\\
\bottomrule
\end{tabular}
\caption{The distant-cue task. Exploratory (session G, S71--S73 and S82); the perfect-gate rows are BOUND, the others ROUTED-BOUND, and session~F's window baseline on other seeds (S70, 330--349) is $0/20$ routed on either memory with $1/20$ and $6/20$ bound.}
\label{tab:far}
\end{table}

\paragraph{The partition does what one channel cannot, and the gate finds it where the cue is local.} Given the partition, every grouped-layout configuration tried with the $N = 256$ model binds, on either write rule. One Hebbian channel does not; one delta channel binds a minority at two streams and none at eight, and the pre-registered version of that comparison is SHOWN on X. A gate that reads three tokens, trained with the memory held slow for 2400 updates and with a channel split when binding stalls, finds the partition at two, four and eight streams without stream labels or restarts, and the split reads the task's key positions.

\paragraph{Discovery is a race, and the recipe's parts each remove one way of losing it.} Slow memory reduces the key splits (12 of 15 failures to 6 of 11 at two streams; the count itself, 7 vs 3, is not shown); the window gate cannot count positions, so position splits are gone; the split breaks the merges that remain at eight streams by giving an idle channel the busy channel's logits, and it is the copy that does it. Two ways of losing remain: at two streams a key split that the copy cannot undo, because, on our reading, there is no idle channel, and that the hinge's key term prevents on a screen (S85, $20/20$ against $15/20$); at a distance, a key split that in the records forms before the register's write at the cue has appeared, which a hinge on the gate's key-dependence during the slow phase prevents ($20/20$ on a screen).

\paragraph{The write rule and the layer.} The delta write with a fixed write strength is not shown worse than the Hebbian one on the grouped layout at two, four and eight streams (one-sided bounds on ten paired seeds each), and better placed for distant cues, where a Hebbian stack needs its first layer routed and a delta stack does not. The partition's claim is about the number of contexts on both rules: one delta channel binds the two-stream header layout $8/10$ and the grouped one $4/10$, by what we read as routes that need no partition, and binds $2/10$ and $0/10$ at four and eight streams.

\section{An independent replication}
\label{sec:replication}

\paragraph{Design.} A separate session (R) was given Revision~8's PDF and Pathway's public BDH repository and nothing else, with the rule that it must not look at this project's code, branches or web presence until its own records were pushed. It implemented the task, the model, the outcomes and the recipes from the text (Phase~1), logging every detail the report did not determine in a file of 44 entries (\texttt{UNSTATED.md}) before running; pre-registered its readings; ran the perfect gate on two seeds and the eight-stream recipe with and without the split on twenty fresh seeds (900--919) at the report's budget (Phase~2); and only then fetched the main line and diffed the two implementations (Phase~3), rerunning ten seeds with one discrepancy at a time set to the main line's choice. Its bar, fixed in advance: ``replicated'' at $\geq 18/20$ bound routed with the split beating the no-split arm at $p < 0.05$; ``partially replicated'' at $\geq 12/20$.

\paragraph{Result.} The perfect gate bound $2/2$. The recipe with the split bound routed $\mathbf{17/20}$ (Wilson $[0.640, 0.948]$); without it $7/20$; every discordant pair favoured the split (10 vs 0, $p = 0.00098$); all 45 of its splits but one copied a channel holding two or more streams onto an empty one; every failure in both arms was a merge of two or three streams, as in ours. The mechanical reading is ``partially replicated'', one run short of the bar. Against the report's $38/40$ the count is Fisher two-sided $p = 0.32$, and pooled over X, L and R the rate is $55/60$ $[0.819, 0.964]$; the referee's assessment (Section~\ref{sec:referee}) is that the bar could not separate the outcomes it was meant to (it fails a true $0.90$ recipe a third of the time) and that the defensible statement is R's own: $17/20$ is consistent with $38/40$, and twenty seeds cannot tell $0.85$ from $0.95$. The two reruns with the main line's probe and with its key positions left the count unchanged on seeds 900--909 ($9/10$, $9/10$), and were not run on the one failure a protocol difference plausibly explains (914, the probe reading of Section~\ref{sec:split}). A fourth phase was proposed, with the recipe as the main line specifies it and a reading calibrated to its size; an amended reading (32 fresh seeds, ``below $0.95$'' at $x \leq 27$, size $0.020$, power $0.54$ at $0.85$) was written after the referee's review but existed in no commit until \texttt{db22857} put it in \texttt{docs/reading/parallel\_sessions\_prompts.md}, after the referee's second pass had found no record of it; R's branch carries the first pre-registered rule, an interval-overlap test pooled with Phase~2 that detects a true rate of $0.70$ about one time in nine, and no Phase~4 record exists.

\paragraph{What the report determined, and what it did not.} R loaded its weights into our model: the gates agree to $7\times10^{-9}$ and the query logits to $7\times10^{-8}$. By R's own labels on its 44 entries, the text stated four outright and three only in part (the gate's input among the four) and left 34 to be guessed, 20 on which it was silent, 12 on which it was ambiguous and 2 carried over from R's own earlier choices, with 3 entries of process (the initialisation scales, the parameter dtype, the evaluation's sampling and the trigger's details among them); in the paths of our code it read (not the CHECK chain, the report printing or the hinge paths) R found nothing that contradicts the report and nothing it had misread. Every difference is protocol, and five could change an outcome: the identity of the split's probe (ours is one fixed set for every seed and arm; R drew one per run), the key positions whose gate mass picks the channels (ours the body's 32 keys; R's included the query's), the perfect gate's schedule (ours a single rate; R's SLOW, which explains its slower oracle), the order of the trigger and the early-stop test at an evaluation step, and the last eligible check. R lists eleven guesses that matched our choices, among them: decay rather than RoPE; that the gate reads the model's own raw embedding; the window gate's width and initialisation, carried over from the recurrent gate; the convolution as one shared width-4 identity filter on the layer input, feeding keys and values; Adam without weight decay and without dropout; the probe's reference at 2400; and the zeroed Adam step counter; the $0.02$ embedding initialisation, which sets the scale of the gate's input, matched as well (R's report), and the map as the argmax of the mean gate at value positions is filed by R as a difference (D-10) that matched in statistic only, since R measured it on the 2048 held-out sequences and the main line on a 256-sequence probe. Appendix~\ref{app:repro} states them, with the entries the referee found missing from this draft's first version (the window's generator, the frozen recurrent weight, the dtype, the perfect gate's model, the tie rules and the noise's estimator) and the three it added (the probe's composition, the outcome's boundary rules, a record format an outside scorer can use).

\paragraph{Caveats.} Isolation rested on self-report. The PDF R was given lay inside a checkout of this repository in its container, and a filename search printed six of that checkout's filenames; R opened none of them and judged the replication not void, and so do we, but a second replication should be isolated by construction (Section~\ref{sec:method}). R used PyTorch 2.14.1 (ours 2.14.0), four workers, and 37 hours of compute, and its records carry no CPU model; bit-for-bit agreement across implementations was neither expected nor sought.

\section{The referee's report and the claim registry}
\label{sec:referee}

\paragraph{What the referee did.} A session with no part in the work (V) was given the repository and asked to review Revision~8 adversarially, changing nothing outside its own directory, which is now \texttt{docs/review\_2026-10/}. It recomputed every count, discordant pair, $p$-value, interval and band of the report's tables from the JSON records under outcome definitions it re-implemented (156 numbers: 140 match, 7 disagree and 6 partly, the 7 all statements corrected in Revision~8.1 and listed below; two later passes over F's reports 3--5, G's 3--4, H's 4--5 and I's 1, 808 rows, found every pre-registered reading re-deriving exactly and 13 disagreements, all in descriptive ranges; the earlier reports and E's batches were audited for when their readings were committed, not re-derived); audited the git history of every main-line claim (no claim or reading edited after a run; the split arm added by a dated amendment before any learned-gate result existed) and of 21 exploratory screens; reproduced one run of each Phase~VI test bit for bit on a 2.10\,GHz Xeon, the model of X's slower host; checked the literature against the project's own notes; and compiled the claim registry of Section~\ref{sec:task}. Its recommendation on the first draft was ``major revision'': the central experiment sound, the manuscript claiming more.

\paragraph{What it found wrong, corrected in Revision~8.1.} Seven statements were false when the first draft was committed: that no single-channel run had been made at eight streams (session~F's premise test was on its branch 82 minutes before the commit, and its counts, $0/10$ at eight streams, bore on the claim); the Muon correction with the two machines the wrong way round (Revision~7's result came from L, was confirmed on X's fresh seeds and reversed on L's); an argument about the convolution that compared an $N = 64$ ceiling model with the $N = 256$ model (the $N = 256$ perfect gate without the convolution binds on every seed recorded); batch~16's hardware; a stale count of single-evaluation binds; a claim that every test had a specification file; and Revision~7's routing ``recorded only on X'' (L's records carry it). Seven statements overstated their evidence and were reworded: ``the reset adds nothing'' (the looks cited showed only that the reset alone is not enough; Section~\ref{sec:split} now has the look that shows it), ``one channel per stream'' (by largest mass; the sets are disjoint), ``no labels'' (no stream labels; the split reads key positions), Gumbel noise ``on the read gate'', ``beat'' for a screen whose verdict was ``promising'', the literature sentence on inferred cues, and the BDH-CQ reconstruction. The referee's check of the corrections found two slips in them (a run count; a phase number) and one older error both had missed (one Hebbian channel \emph{had} been run at four streams: $0/20$), corrected in a second pass.

\paragraph{What it asked for, and what this revision does with it.} Scope (done: the title and Section~\ref{sec:limits}); the copy-only arm as a pre-registered claim (Section~\ref{sec:split}; the copy against the reset is pre-registered and shown on X, the copy alone against the full split is next); the split's dependence on key positions (stated everywhere it is claimed; a trigger that reads all positions is not yet screened); an external baseline family (one delta channel and the delta write with its decay and write strength fixed are now run at two, four and eight streams on the grouped layout, and Gated DeltaNet's learned forget gate on the eight-key header layout; no external model, a Gated DeltaNet, Mamba or attention model, has been trained on the task; a residual-stream router was run and failed; Mixture-of-Memories' balance loss was run); paired-difference intervals and a stated error rate (Section~\ref{sec:main}, Section~\ref{sec:task}); routing by overlap (reported); the literature's remaining overstatements (Section~\ref{sec:memory}); provenance (commits pinned; L's local start times noted; L's Phase~III--IV records committed, Appendix~\ref{app:prov}). Two requests are open: the claim list's $q$-values are in the registry and not yet in the tables here, and BDH-CQ and Triadic Linear Attention have still not been read in full.

\paragraph{The registry's two findings about the records.} L's result files for Phases~III--IV were not in the repository: the commit that was meant to hold them (5 October) held byte-identical copies of X's files, naming X's processor, so every L verdict from those phases was known only from text written later. The user recovered twelve of the thirteen files from the laptop's working copies and reran the thirteenth; all 70 L rows the registry can now check against a file agree with it (38 verdicts, of which 17 had been second-hand, and the 32 bands, gates and readings recounted from the per-run records), eleven L tests gained a printed $p$, and L's family under correction improved (27 of 32 positive verdicts survive Benjamini--Hochberg). The main line's split probe is unbalanced across streams (Section~\ref{sec:split}).

\paragraph{The sessions' reports, after the referee.} Statements the sessions made beyond their pre-fixed readings are quoted in this revision as descriptions, not results: among them, that S54's slow delta oracle was a learned-$\beta$ artefact, that the register's write at the cue is what separates a discovering run, that the stack's stream tag lets the gate stop routing, and that the check interval and the first eligible check act through the flat start. Each is marked where it appears; the referee's check of this draft's first version listed the places where they were not, and those are now marked.

\paragraph{What it found in this revision's first draft.} The referee checked the first draft of this revision (commit \texttt{9bfa85a}) against the records, the docstrings and the code (\texttt{docs/review\_2026-10/revision\_9\_check.md}, 400 rows): every count, discordant pair, $p$-value and interval in its tables reproduces, every reading it quotes was committed before its first run, and its text repeated Revision~8's failure in kind. False when committed, against records already pushed: that the main line's records store neither the probe's reading nor the per-stream masses at each check (they store both; Section~\ref{sec:split}); that session~G was still screening three devices whose result, $20/20$, had been recorded fifteen minutes before the commit (Section~\ref{sec:far}); that session~H's S85 was running (committed, not started); that session~R's Phase~4 reading had been corrected (the amended reading existed in no commit until \texttt{db22857}; R's branch carries the first rule); that the referee's recomputation found ``all'' numbers to reproduce (140 of 156) and had re-derived every reading of every screen (the later reports; the earlier ones were audited for timing); that 90 L rows agree with the registry (70); that a run was reproduced on a third CPU model (on X's own slower model); and that R's model agrees with ours to $10^{-8}$ (the logits to $7\times10^{-8}$). Contradicted by the draft's own tables or by the record: that the reset-only arm bound no more often than no trigger in every look (S61: $4/10$ against $3/10$); that all 83 of S61's firings copied a row (57; the resets copy nothing); that a learned $\beta$ collapses (it rose in two screens); that six delta splits fired before the map had merged (it had; five of the six copied the wrong channel); and three ranges stated for every run that held for the median run. Overstated: the delta write's ``parity'' and ``interchangeable'', the S61 pattern in a new place, now ``not shown worse''; ``settled'' for screens; ``supported'' citing a screen whose reading did not apply. Unmarked descriptions stated as findings, and readings attributed to the wrong arm or quoted without their docstring's qualifier, in some fifteen places. The title still named the general problem. Every item is addressed in this version, which the referee's second pass checked again (\texttt{rev9/resolution\_*.md}; it found the second version's own slips, among them a miscount of R's labels and the uncommitted Phase~4 reading, corrected in turn), and Appendix~\ref{app:repro} gained the six entries the check found missing.

\section{Retrospective}
\label{sec:retro}

\subsection{Revision~8's claims}

\paragraph{``Eight streams bind without labels or restarts'' (the main result).} \textbf{Stands}, and was replicated from the text alone ($17/20$; Section~\ref{sec:replication}). Qualified: without \emph{stream} labels; the split reads key positions; ``one channel per stream'' by largest mass, the sets disjoint.

\paragraph{``The copy, not the reset.''} \textbf{Pre-registered and shown on one machine} (C1, 7 vs 0; C2 holds), pooled pending L. The first draft's ``the reset adds nothing'' was not shown by the looks it cited; session~H's S61 holds it as a one-pair bound on ten seeds (the copy with noise $9/10$, the exact copy $9/10$, the full split $9/10$), and the pre-registered version of that comparison is next.

\paragraph{``The Hebbian write stays on the main line.''} \textbf{Withdrawn.} With the write strength fixed, the delta write is not shown worse than the Hebbian one at two, four and eight streams on screens and binds sooner at two and eight (Section~\ref{sec:memory}); the stall Revision~8 reported came from a screen with a learned write strength (and a different split and seeds; the attribution is post hoc). A pre-registered test is specified.

\paragraph{``The partition's claim is about the number of contexts.''} \textbf{Supported on X} by the pre-registered C3 ($0/10$ against the recipe's $10/10$); session~F's S68 had $0/10$ at eight streams too ($2/10$ at four), but its own reading of that name did not apply.

\paragraph{``If the stream-tagged keys at the third layer hold up, the context is reaching depth.''} \textbf{They did not} ($1$ of 16 bound runs; Section~\ref{sec:memory}); on the randomised header layout the context reaches depth in the delta solutions that bind without routing, and on the grouped two-stream layout it reaches depth at key positions in routed runs as well (Section~\ref{sec:memory}).

\paragraph{``The next steps are a probe, a register and a per-layer residual gate.''} \textbf{Done.} The probe found the cue absent from the Hebbian stack; the residual gate alone bound nothing on Hebbian memory and routed nothing on delta; the register carries the cue once it also gates the first layer, which the first design did not do; that this design choice is what sank the Hebbian register is our interpretation, across two tasks and two gate inputs (Section~\ref{sec:far}).

\paragraph{``Scale-free'' at four streams (batch~19).} \textbf{Stands as a screen}; the referee notes that its one-pair rule on ten seeds reads ``scale-free'' $39\%$ of the time when the larger model's rate is $0.7$ against $0.9$, and that ``depth'' is applications of one shared layer. Eight streams at the larger sizes (S66, S67) are running.

\subsection{Revision~7's claims}

\paragraph{``At eight streams the gate forgets the context'' (subtitle).} \textbf{Superseded.} True of the recurrent gate; the window gate carries the context at every write position by construction and binds $38/40$.

\paragraph{Revision~7's recipe (slow memory and a hinge) binds without restarts at two streams and at four with sixteen channels.} \textbf{Stands}, and is replaced by a recipe without the hinge, which is within eight points of it at two streams ($[-0.079, +0.121]$), beats its early-window form at four (G1) and beats it at eight (G2).

\paragraph{``The gate needs only one token of memory'' (Section~11 of Revision~7).} \textbf{Corrected.} One at the read, two at the write; the write-position decoders were not measured in Revision~7's screens S38--S41 (batch~16, S47, measured them).

\paragraph{``Under Muon the recipe binds four streams sooner.''} \textbf{Not shown} on disjoint seeds pooled (12 vs 8, $p = 0.25$), and the machines disagreed: confirmed on X's fresh seeds (7 vs 1), reversed on L's (5 vs 7), where Revision~7's result had been shown. Revision~8's first draft called it a seed effect, which the data do not establish; the referee's heterogeneity test between the machines ($p = 0.070$) is the more interesting figure, and a seed-block swap between X and L would separate processor from seeds.

\paragraph{``A cap on the gate's recurrence is being screened.''} \textbf{Resolved: not.} No capped arm bound ($0/10$, $0/10$, $0/10$ and $0/5$); the cap acted from update 1.

\paragraph{The row copy breaks merges (Section~10 of Revision~7).} \textbf{Supported} with the right outcome and a reset-only control in four exploratory looks and, on X, by the pre-registered C1 (the copy with its reset against the reset alone, 7 vs 0); that the runs it rescues are merges is printed in their failure classes, not claimed.

\subsection{Errors made during Phase~VI}
\label{sec:errors}

\paragraph{Statements that were false when committed.} Listed in Section~\ref{sec:referee}. The common cause was writing from a summary of the branches rather than from the branches: the premise test and the copy-only screen had both been pushed before the draft was committed, and neither was checked. The fact-check that preceded the draft was done by the author of the draft. This revision's first draft repeated the failure in kind: the referee's check of it found nine statements false against records already pushed, among them that the records do not store the probe's reading, that session~G was still screening devices whose result was fifteen minutes old, that a screen was running that had not started, that a replication rule had been corrected that had not, and that the recomputation found ``all'' numbers reproducing; they are listed in Section~\ref{sec:referee} and corrected here. The cause was the same, and the remedy is the one the referee applied: every statement about a branch checked against that branch at its head, by a script, before the commit.

\paragraph{Predictions that the data overturned.} We predicted that a single delta channel could not bind conflicting streams (it binds the two-stream header layout $8/10$); that Gumbel noise on the gate's logits could replace the split (on both gates it collapsed routing; on the write gate it bound $0/10$ at eight streams); that a balance loss would attack merges (it changed the key- and value-position maps instead); that a learned forget gate would widen the gap between the oracle and one channel (it shrank it, on the eight-key header layout); that the register at layers 2--3 would carry the cue on Hebbian memory (it needed layer~1, on our reading of S71); that the delta recipe would stay behind at eight streams (it keeps them $9/10$ against $10/10$ with the write strength fixed); that a single Hebbian channel's stream tag, if any, would be linearly readable (it is not); that a dense loss would need the partition at two sources (one channel self-tags); that a gate held slow at the start would let the register's write win the race against the key split (it binds $0/20$; on our reading it slows the write too) and would prevent the two-stream key split (it has four key failures of its own and $14/20$ against $15/20$; its three onsets after update 0 come later than the baseline's seven, on two shared seeds); and, in the audit, that the window gate with slow memory would reach $\geq 36/40$ on seeds 160--199 with no key failures (it reached $31/40$ there, and $33$ and $36$ of 40 on the main test's seeds, with key failures reduced).

\paragraph{A design choice that cost a round.} Revision~8's register specification kept the first layer on the window gate on the reasoning that its residual input is the embedding and carries no cue; but the register is built from embeddings and could have fed the first layer's gate, and on Hebbian memory the first layer is the one whose routing matters. One screen (S71) settled it; it could have been part of the first round.

\paragraph{A replication bar that could not discriminate.} Session~R's ``$\geq 18/20$'' fails a true $0.90$ recipe one time in three, and the first version of its Phase~4 reading, written by us, pooled the seeds already seen and would have detected only rates near $0.70$. The referee's review corrected both in principle; an amended Phase~4 reading was written after it but existed in no commit until \texttt{db22857}, after the referee's second pass had found no record of it, and R's branch carries the first rule. Section~\ref{sec:method} gives the rule a replication should use.

\paragraph{A probe that can be wrong, and is unbalanced.} The split's 64-sequence probe can read a plateau as bound (R's seed 914) and weights the streams unequally, in the same way for every seed (Section~\ref{sec:split}). Both could have been read from the project's own records, which store the probe's reading at every check beside the evaluation's; nobody read the two against each other until the referee did (6 of 353 checks; Section~\ref{sec:split}).

\paragraph{A distant-cue task that exceeded its oracle.} Session~F's first randomised header layout did not bind under the perfect gate within its budget; the lighter version with four keys per stream in two blocks did, and is the task of Section~\ref{sec:far}.

\paragraph{A classifier at its own threshold.} Three batch-18 runs sit within $10^{-3}$ of $\eta^2 = 0.5$ by key and were labelled STREAM-PARTIAL; the key-first classifier relabels the one at or above $0.5$ KEY (Section~\ref{sec:task}). The reading that depended on them is withdrawn.

\paragraph{Records that were not L's.} A commit of ``machine L's test results'' held copies of X's files for 20 tests; the real files were recovered from the laptop and one test rerun (Section~\ref{sec:referee}).

\paragraph{Branch discipline.} Session~I committed its first files to the main line instead of its branch; they were moved, byte-identical, and the main line cleaned. Two sessions shifted into the same seed block (1300--1359) after collisions with earlier blocks; the next round's blocks were checked against every branch before assignment.

\paragraph{Attribution of a number.} An earlier working note put Pathway's $97.4\%$ Sudoku figure beside TRM's $87.4\%$ as if on one protocol; the figure is from Pathway's blog, on internal data, and the protocols differ.

\paragraph{Process.} Container restarts cost in-flight runs on X (the test's first start), on F, G and H (while sessions were idle with only detached processes running) and on E in several batches; a detached launcher with a watcher that resumes after a restart (\texttt{longrun.py}) now runs every main-line test, and every resumed segment begins by reproducing a recorded run bit for bit. Batch~18's Muon reference did not reproduce across hosts and was rerun fresh, as the rules require.

\section{Methodological findings}
\label{sec:method}

Revisions~5--7 recorded sixteen lessons and Revision~8 five; this revision adds seven, most of them learned from being reviewed.

\paragraph{Disjoint seeds make the pooled test a test.} From the Muon-recipe test on, each machine has its own seeds and the pooled McNemar is the primary claim. The window-gate test's G1 is the case this was for: not shown on X alone, shown on L and pooled.

\paragraph{Audit the records before the next screen.} The window gate had been screened, had won, and had been dropped for a reason unrelated to the obstacle at hand; one reading of the records found it. A fixed prediction before the rerun made the audit falsifiable.

\paragraph{Control for the bundled action.} A split that copies and resets needed a reset-only arm and a copy-only arm; so does every intervention with two parts.

\paragraph{Score routing at the write positions.} A stream's channel at key positions can differ from its channel at value positions in a bound run (session~H: 34 of the balance-loss arm's 36 bound-routed runs, and two to five of the 37--39 bound-routed runs in each of S58's other full arms); the writes are at the values.

\paragraph{Borrowed devices need the model's read path.} Gumbel noise is the device of a router whose reads are dense; on our routed read it collapsed routing. A balance loss balances whatever is cheapest to balance.

\paragraph{Have the draft checked by someone who did not write it.} The first draft's fact-check was done by its author and missed seven false statements, two of them about results that were already on the branches. An adversarial referee with the records found them in a day, recomputed the numbers of the tables, and found the numbers right; the text was the problem. The same happened to this revision's first draft, checked by the same referee against the branches at their heads (Section~\ref{sec:referee}): the numbers right, nine statements false. The referee now reads every main-line specification before its first run and every draft before it is merged, and a statement about a branch is checked against that branch by a script before the commit.

\paragraph{Fix the name of the outcome, not only the rule.} A reading named ``discovered'' on a screen that also prints an outcome named DISCOVERED came out ``does not apply'' under its rule and would have applied under the outcome of the same name (4 vs 0 against 5 vs 0). Nothing was misread, but the next reader could have. A reading's outcome should be the only outcome that bears its name.

\paragraph{Report descriptions as descriptions.} The sessions' reports, and ours, state mechanisms beside readings: that a learned write strength caused a stall, that the flat start wastes splits, that a write at the cue is what separates discovering runs. The referee found every such sentence to lack a pre-fixed rule, and some to rest on a count that no classification reproduces. They are kept, marked, and none is a result.

\paragraph{A replication needs a calibrated bar.} ``Replicated'' at $\geq 18/20$ fails a true $0.90$ recipe one time in three; an interval-overlap rule on pooled seeds detects a true rate of $0.70$ about one time in nine and rarely any rate above $0.65$. The rule for a replication of a rate $r_0$ on $n$ fresh seeds is a one-sided exact test against $r_0$ with its size and its power at the alternative that matters stated in advance ($x \leq 46/52$ against $0.95$ has size $0.045$ and power $0.81$ at $0.85$), beside a Fisher comparison with the original count and a paired test of the mechanism's arms on the same seeds; isolation by construction, in a container that never holds the repository; records in a format the original's outcome code can score.

\paragraph{Read what the trigger saw.} The records keep every check's probe reading beside the evaluation's, and nobody read the two against each other, so a probe that read a plateau as bound, and a probe that under-weights some streams, were found by an outsider. A test should print them side by side, with the count of checks the probe blocked, as \texttt{test\_delta\_recipe} now does.

\paragraph{Branches and seeds are part of the protocol.} A session that pushes to the wrong branch, a seed block that collides with another session's, and a commit of ``L's results'' that holds X's files are all provenance errors that cost nothing to prevent and hours to repair; each happened once.

\section{Limitations}
\label{sec:limits}

\paragraph{Scope of the main claim.} The result is that a gate trained only on the task loss finds a partition of the memory on the \emph{grouped layout} of a synthetic binding task, where every write's context cue is within three tokens, at two, four and eight streams, with a split aimed by gate mass at the task's key positions. The title and subtitle say so. ``Context-conditional binding'' in general, natural data and the mechanism's role in a language model are not shown; the distant-cue task has a solution without stream labels on screens ($20/20$ with a hinge that reads key identities during the first 2400 updates) and no pre-registered test of it yet.

\paragraph{The recipe needs task structure.} The split reads the key positions, and the hinge's key term, which on screens prevents the key split at two streams and for the register at a distance, reads key positions and key identities; a trigger that reads all positions has not been screened. The records say what one would have done: every check of every split run stores the channel pair an all-position mass would have chosen beside the key-position pair, and at the 107 firings of the 50 recorded split runs at eight streams (\texttt{window\_gate} X and L, \texttt{split\_copy} X) the all-position pair was a labelled target, a channel holding two or more streams and an empty one, at 97 (21 of 22, 46 of 47 and 30 of 38), the key-position pair at all 107; the two pairs coincide at 46. That is our count, not a reading, and a run that split differently would have gone differently, so it is a case for a screen, not a result. A language model has no marked keys; the analogue is untested.

\paragraph{The model.} Three layers with shared weights (two to six screened), one head, $d = 32$ (to 128 screened), $N = 256$ (to 1024 screened), sequences of at most 99 tokens on the main line and 480 in session~I's pilot, a fixed memory half-life of about 14 tokens on the main line, which is the horizon of everything the memory can carry; a language model's would need to be thousands, and the one device screened against it, a learned forget gate, was screened once at a fixed length: it raised one delta channel's binding ($9/10$ and $8/10$ against $3/10$) and not a Hebbian one's, and its pre-fixed reading on the horizon read ``it does not'' (Section~\ref{sec:memory}).

\paragraph{Two streams.} The hardest configuration relative to its size; every arm is at $65$--$69/80$, the one repair device that works at eight streams did not work here (S78, ``neither device''), and the prevention that works on a screen (the key term, $20/20$) reads key identities.

\paragraph{Screens.} Exploratory sessions use 10--40 runs per arm and lenient rules; their readings are labelled, their descriptions are descriptions, and nothing from them is a claim until a main-line test repeats it. Several screens in this revision have both arms at the ceiling and can detect only large differences.

\paragraph{Replication.} One, from the text, by one implementer, on twenty seeds, isolated by self-report; it agrees in direction and size and is one run short of its own bar.

\section{Distance to a language model}
\label{sec:lm}

\paragraph{What Phase~VI adds to the case.} The gate is a convolution: parallel, cheap, the form every modern linear-attention model already uses. The recipe is two learning-rate groups and a plateau rule. It finds an eight-way partition without stream labels, its constants held, on a ten-seed screen, to four times the width and twice the depth at four streams, and its write rule can be the delta rule that the linear-attention literature already uses.

\paragraph{A dense loss (session I).} Revision~8 conjectured that an end-of-sequence loss gives too little pressure to discover a distant cue, and that language, with a loss at every token, supplies more. Session~I built a task to ask it: $S$ sources, each with its own table of three successors per word drawn fresh for every sequence (so the sets must be learned in context), interleaved in blocks of random length and random order with a source token at each block's start, and a next-token loss at every word with a known same-source predecessor (the task family of \cite{bietti, edelman}, cited for it; neither paper has been read in full). Its pilot ran the perfect gate (one channel per source), one channel and a single-source control at two sources, 12{,}000 updates, four seeds, and could not make the task usable by its own rule, because \emph{one channel learns to carry the source}: in three of four runs of the shorter-block configuration a step transition raised the set accuracy by $0.09$--$0.22$ and took a linear probe of the source, from the residual entering layer~3, from $0.54$--$0.58$ to $0.96$--$0.98$, closing $44$--$83\%$ of each seed's gap to the perfect gate with no gate and no labels; the perfect gate's own stack never carried the source, having no need to. The sparse-loss stacks of Section~\ref{sec:far} never did this. So the dense loss does put pressure on the stack to carry the cue, which is the conjecture's first direct evidence and is descriptive; whether a gate that reads such a stack then partitions the memory, and whether four sources are too many for one channel to tag, are the questions, and four sources need sequences of about 770 tokens, which is a GPU's job. Session~I's GPU package (its branch; a device-agnostic runner for the project's models and tasks with a CPU check against a recorded run and a CUDA check at a stated tolerance) estimates that screen, from a CPU profile, at 9--13 hours on an RTX 3070 Laptop or 14--21 on a GTX 1660 SUPER.

\paragraph{What stands in the way,} in order:
\begin{enumerate}
\item \emph{Discovery of a distant cue without task structure.} The architecture exists and, with the key hinge, is found $20/20$ (Section~\ref{sec:far}); without it the write at the cue is found on a fifth of seeds in S72 and S73 and on 2 of 20 fresh seeds in S82, and a key split, formed in the records before any write at the cue, ends fourteen of the sixteen others on each memory. The hinge reads key identities; a dense loss may change the race without it.
\item \emph{The split's nuisance variable.} The key positions come from the task; a trigger on all positions is the next screen.
\item \emph{Memory horizon.} A fixed decay; a learned forget gate raises one delta channel's binding at a fixed length and not a Hebbian one's, its reading on the horizon itself read ``it does not'', and the delta write is now a candidate for the main line.
\item \emph{Baselines.} One channel of either write rule is run at two, four and eight streams; Gated DeltaNet's write is the delta rule; Mixture-of-Memories' balance loss and Raven's noise were run; a residual-stream router was run. The comparison that matters is whether their routers discover \emph{contexts} rather than token types, which none of their papers tests, and which the register result now makes askable.
\end{enumerate}

\paragraph{A staged path.} (i) On CPUs: the memory-rule test and the register test on the main line; the key-split devices; eight streams at the larger widths. (ii) On the user's GPUs, with session~I's package: the dense-loss task at four sources; the recall benchmarks \cite{zoology}; a replication on a second machine class. (iii) Small language models, asking one question: do the channels specialize by context or by token type.

\section{Conclusion and next steps}

A gate trained only on the task loss finds the partition of a Hebbian working memory that context-conditional binding needs, on the grouped layout of a synthetic task, at two, four and eight streams: a window gate, slow memory for 2400 updates, and a channel split when binding stalls, in 38 of 40 runs on two machines with disjoint seeds, where the previous recipe bound none, and in 17 of 20 in a replication from the text alone. On one machine so far, pre-registered, the copy with its reset beats the reset alone, the reset alone is within the bound of no trigger, and one delta channel binds none of ten seeds where the recipe binds all ten; the pooled readings wait for the other machine. The write rule is open: with its write strength fixed, the delta rule is not shown worse than the Hebbian one on screens at every size tried. At a distance, a register that overwrites on the cue and gates all three layers carries it when told where to split, is found label-free on a fifth of seeds (2 of 20 on the fresh seeds of the device screen), and is found on every seed of a screen when a hinge keeps the gate off key identity during the slow phase, which on a screen also prevents the two-stream key split that the copy did not repair; on Hebbian memory the first layer is the one that must be routed. Under a dense loss one channel learns to carry the cue by itself. A referee recomputed the numbers of Revision~8 and of this revision's first draft from the records and found them to reproduce; the statements it found wrong in each draft are corrected and listed.

Next:
\begin{itemize}
\item L's half of \texttt{test\_split\_copy}, then \texttt{test\_delta\_recipe} on both machines (specified, reviewed and amended once), then the hinged-register test on the randomised header layout (session~G's proposal: the hinged register against the baseline, the plain register and the hinge without the register; 20 seeds per machine; a rate claim of $\geq 18/20$).
\item Session~H's proposed test of the key term at two streams (the window recipe with and without it, disjoint seeds, the key-first classifier); session~F's trigger on delta channels (S81); batch~19's eight streams at the larger sizes; session~R's Phase~4 under the amended reading, now committed.
\item The dense-loss task at four sources on a GPU, once session~I's package passes its CUDA check.
\item Revision~10 after those; BDH-CQ and Triadic Linear Attention read in full before it.
\end{itemize}

\appendix
\section{Provenance}
\label{app:prov}

\begin{table}[H]
\centering
\footnotesize
\begin{tabular}{@{}>{\raggedright\arraybackslash}p{0.15\linewidth}>{\raggedright\arraybackslash}p{0.4\linewidth}>{\raggedright\arraybackslash}p{0.28\linewidth}>{\raggedright\arraybackslash}p{0.09\linewidth}@{}}
\toprule
test & X: spec / test commit / result commit & L: results file (commit recorded) & runs, X; L\\
\midrule
\texttt{muon\_recipe} & --- / \texttt{74f5907} / \texttt{3e2d171}; pooled \texttt{2a6ba3d} & \texttt{L/muon\_recipe\_results.json} (\texttt{74f5907}) & 156; 166\\
\texttt{window\_gate} & spec \texttt{9a769c7}, amended \texttt{d05f0ca} / \texttt{2f98f06}, Part~D \texttt{886a668} / \texttt{628eb7c} & \texttt{L/window\_gate\_results.json} (\texttt{886a668}; recorded \texttt{1cd20a1}) & 231; 231\\
\texttt{split\_copy} & spec \texttt{fb3acb7} / \texttt{1f52cd5} / \texttt{592c574} (X; the classifier \texttt{ba8901a}, the launcher \texttt{56c44a7}) & not run & 54; ---\\
\texttt{delta\_recipe} & spec \texttt{05256c5}, Amendment~1 \texttt{32c121b} / test \texttt{a581027} / not run & not run & ---; ---\\
\bottomrule
\end{tabular}
\caption{Phase~VI main-line tests on branch \texttt{claude/bdh-growth-hebbian-inference-w90069}. L's results files for Phases~III--IV are in \texttt{results/L/} from \texttt{9a53302} (twelve files, recovered from the laptop's working copies) and \texttt{9c53e16} (\texttt{channel\_binding}, rerun on L on 10 October with the test's code unchanged); the Phase~V files at \texttt{00ac0f2}, \texttt{d0e6c85}, the scale-axes file at \texttt{42af095}. X: Intel Xeon at 2.80\,GHz for the first two tests and 2.10\,GHz for \texttt{split\_copy}; L: Intel i7-12650H; PyTorch 2.14.0 (L's rerun 2.14.1). L's recorded start times are local time (UTC$-$6); X's are UTC.}
\label{tab:prov}
\end{table}

\begin{table}[H]
\centering
\footnotesize
\setlength{\tabcolsep}{3pt}
\begin{tabular}{@{}>{\raggedright\arraybackslash}p{0.11\linewidth}>{\raggedright\arraybackslash}p{0.52\linewidth}>{\raggedright\arraybackslash}p{0.16\linewidth}>{\raggedright\arraybackslash}p{0.15\linewidth}@{}}
\toprule
batch or session & screens (report) & commit & runs\\
\midrule
15--18 (E) & S41--S52 (Revision~8) & \texttt{90c3a28}, \texttt{93bdfe6}, \texttt{2531480}, \texttt{7ef4e15} & 379\\
19 (E, running) & S64 width, S65 depth; S66 eight streams and S67 capacity under way & \texttt{8cff21e} & 118 at \texttt{5dfc2b5}\\
H & S58, S58b, S60 (reports 1--2); S61 copy without reset, S62 query gate, S63 hardness (report 3, \texttt{7e9197f}); S78 two streams (report 4, \texttt{72f0bbc}); S79 constants (report 5, \texttt{df57087}); S84 guard and evaluation trigger (report 6, \texttt{37c1fbf}); S85 prevention at two streams (report 7, \texttt{1e09ecd}) & \texttt{1e09ecd} & 318; 300; 82; 62; 52; 82\\
F & S53, S54, S59 (reports 1--2); S68 premise, S69 Hebbian decay, S70 lite task (report 3, \texttt{9a6c2bf}); S77 probe (report 4, \texttt{22ec912}); S76 eight streams on delta (report 5, \texttt{5b9bc2a}); S80 two and four streams (report 6, \texttt{ed76eea}); S81 running & \texttt{ed76eea} & 322; 102; 30; 40; 48\\
G & S55 probe (report 1, \texttt{6cd36ff}); S56 register, S57 latch (report 2, \texttt{427db36}); S71 layers (report 3, \texttt{35de6db}); S73, S72 register at all layers (report 4, \texttt{93422cf}); S82 devices against the key split (report 5, \texttt{8973d48}; spec \texttt{83699a2}); S83 the hinge's inputs, the control and the window (spec \texttt{b2f9907}; running) & \texttt{8973d48} & 22; 139; 50; 150; 136\\
I & S74 pilot, three configurations (report 1, \texttt{a2588da}); the GPU package (in progress) & \texttt{a2588da} & pilot\\
R & Phase 2 (\texttt{de8df0f}), Phase 3 diff and reruns (\texttt{00990ea}, \texttt{100dced}); Phase 4 rule (\texttt{464be8d}), not started & \texttt{464be8d} & 42; 20\\
V & review (\texttt{d2bb959}); 8.1 check (\texttt{0fdb363}); registry (\texttt{ebb9325}); replication review (\texttt{01bee1f}); session recomputes (\texttt{4bf85e9}); L addendum (\texttt{0e6c189}); delta spec reviews (\texttt{fad1489}, \texttt{16b04f9}); the test check (\texttt{6f2788d}); this revision's checks (\texttt{e8fbb25}, \texttt{66334b8}) & \texttt{66334b8} & 2 reproductions\\
\bottomrule
\end{tabular}
\caption{Exploratory work and the review in Phase~VI, one thread per run, each segment opened by a bit-for-bit reproduction of a recorded run; one commit per report as listed, the branch's commit column being the last report's. Branches \texttt{claude/outside-ideas}, \texttt{claude/explore-F/G/H/I}, \texttt{claude/replicate-R}, \texttt{claude/review-V}; the review is merged into the main line under \texttt{docs/review\_2026-10/}. The audit is \texttt{docs/audit\_2026-10.md} (\texttt{b22da04}); the reading notes and synthesis are \texttt{docs/reading/} (\texttt{ad8a92e}); the sessions' prompts are \texttt{docs/reading/parallel\_sessions\_prompts.md}.}
\label{tab:batches}
\end{table}

\section{What a replication needs: the protocol as this revision states it}
\label{app:repro}

The items a reader must have to score runs as the main line does, compiled from session~R's \texttt{UNSTATED.md} (44 entries), its \texttt{DIFF.md} (17) and the referee's additions. Where the report had been silent, the main line's choice is given.

\paragraph{Model.} Pathway's BDH forward pass in float32 with $x = \mathrm{LN}(\mathrm{emb})$ and, per layer, the sparse code $\sigma = \mathrm{ReLU}(xE)$, $y = \mathrm{LN}(\mathrm{attn}(\sigma, \sigma, x))$, $y_\sigma = \mathrm{ReLU}(yE_v)$, $z = \mathrm{LN}((\sigma \odot y_\sigma)D)$ and $x \gets \mathrm{LN}(x + z)$, every LayerNorm without affine parameters, and logits $x W_{\mathrm{head}}$ after the last layer; one head, $N = 256$, $d = 32$, three applications of one shared layer; scalar decay $0.95^{t-u}$ in place of RoPE; strictly causal ($u < t$); no dropout; Pathway's initialisation (encoder, decoder, head and embedding at $\mathcal{N}(0, 0.02^2)$), which sets the scale of the gate's input; channel mixing as the score $(\sigma_t \cdot \sigma_u)(g_t \cdot g_u)\,0.95^{t-u}$ on the sparse codes, with one gate per token for read and write, the same at every layer; the gate reads the model's own raw embedding (the same table), through a width-3 depthwise causal convolution (weights uniform in $\pm 1/\sqrt{3}$ from \texttt{Generator(seed + 31337)}, no bias, left-padded with zeros), a hidden layer of 32 units with $\tanh$, $W_{\mathrm{in}}, W_g \sim \mathcal{N}(0, 0.1^2)$, no biases; the model is built as the recurrent-gate model and converted, so its recurrent weight $W_h$ is drawn (which fixes the draw order of the parameters after it) and stays frozen, unused and in no optimizer group; the ``conv'' of the four- and eight-stream configurations is one shared depthwise causal filter of width 4 on the $d$-dimensional layer input, initialised to the identity, whose output feeds the keys (equal to the queries, through $E$) and the values; the perfect gate's model is \texttt{test\_stream\_recipe}'s \texttt{ceiling8k16} ($k = 16$, stream $s$ to channel $s$, the context-token list padded with $16 - S$ ids that never occur so that channels $S$ to 15 are unused), with the curriculum off.

\paragraph{Task.} 28 tokens at eight streams and four keys: 16 value tokens, $P$ key tokens, $S$ context tokens, and no query token (the sequence of 99 tokens ends with the query's value; the model reads the first 98 and predicts the 99th from the query's KEY position); per sequence and key, $S$ distinct values without replacement, assigned to streams in random order; the query's stream and key uniform and independent; the loss is cross-entropy over the full vocabulary at the query's KEY position; accuracy is argmax over the full vocabulary; 2048 held-out sequences from a fixed generator (\texttt{Generator(424242)}), the same for every run and arm; \texttt{torch.manual\_seed(seed)} before initialisation, the training batches from \texttt{Generator(seed + 10000)}.

\paragraph{Training and evaluation.} Adam with default betas, no weight decay, batch 32, constant rates, no warm-up, no clipping; SLOW's groups and its switch after the evaluation at 2400; evaluation every 1200 updates and at the budget; early stop after three consecutive evaluations $\geq 0.95$; at a shared step the order is evaluation, statistics, the path's own check, the learning-rate switch, the split trigger, the early-stop test; the perfect gate trains at the single rate $10^{-3}$ without SLOW; the final per-stream accuracy and map come from a final evaluation of the final model.

\paragraph{Outcomes.} Bound: the first evaluation $\geq 0.95$ that holds to the end (a run can bind on its last evaluation); the stream-to-channel map: for each stream, the argmax over channels of the mean gate at the body's value positions on a 256-sequence probe (\texttt{Generator(seed + 99000)}), ties on unrounded values to the lowest channel index (\texttt{torch.argmax}); bound routed: bound, the map one-to-one, every stream's accuracy $\geq 0.9$ on the held-out queries of that stream; failure labels as Section~\ref{sec:task}.

\paragraph{The split.} The probe: 64 sequences from \texttt{BindTask.make\_batch(64, Generator(12345))}, the same for every seed and arm, scored on their queries (5, 5, 6, 8, 10, 12, 5 and 13 per stream at eight streams); read at 2400 as the first reference and at every check from 4800 in steps of 2400 up to the budget minus 2400; fire if the probe's accuracy is below $0.95$ and rose by less than $0.02$ since the previous check; $c^*$ and $c_0$ the argmax and argmin of the mean read-gate mass over the probe's 32 body key positions (not the query's), ties to the lowest index; $W_g[c_0] \gets W_g[c^*]$ (the main line's Part~D form adds $0.1\,\sigma$ noise to both rows, $\sigma$ the copied row's standard deviation by \texttt{torch.std} with Bessel's correction, two draws of the row's shape from \texttt{Generator(12\_000\_000 + 1000*seed + n\_splits)}, the first added to $c^*$'s row and the second to $c_0$'s, and zeroes $W_g$'s Adam state including its step counter); a split blocks further splits for 4800 updates; at most three; no guard against a target that carries key mass (none is needed at $k = 16$, $S = 8$).

\paragraph{Records.} Per run: the curve at every evaluation, the per-stream gate mass at value positions at every evaluation and the per-stream accuracy at the end, each check's probe reading beside the evaluation's reading at the same step, each split's update, $c^*$, $c_0$ and the map before it, the CPU, the code commit and the start time with its time zone. The main line's records hold all of these except the zone (its start times are local to the machine, without one), and a replication should record it.

\paragraph{Statistics.} Exact one-sided McNemar on the discordant pairs, $p = P[\mathrm{Bin}(b + c, \tfrac12) \geq b]$; Wilson 95\% intervals with $z = 1.96$; for a replication of a rate on $n$ fresh seeds, a one-sided exact test against the reported rate with its size and power stated in advance, a Fisher comparison with the original count, and the mechanism's arms paired on the same seeds.

\begin{thebibliography}{99}
\small
\bibitem{bdh} A.~Kosowski, P.~Uzna\'nski, J.~Chorowski, Z.~Stamirowska, M.~Bartoszkiewicz. \emph{The Dragon Hatchling: The Missing Link between the Transformer and Models of the Brain}. arXiv:2509.26507, 2025.
\bibitem{bdhcq} B.~Engdahl, A.~Kosowski, J.~Chorowski, Z.~Stamirowska, P.~Uzna\'nski, J.~Jiang, R.~Phadke, R.~Kinas, R.~Zhong. \emph{BDH-CQ: In-Context Learning with Recurrent Latent Reasoning}. arXiv:2608.09888, 2026.
\bibitem{tla} O.~Sieberling et al. \emph{Triadic Linear Attention}. arXiv:2609.36529, 2026.
\bibitem{mamba} A.~Gu, T.~Dao. \emph{Mamba: Linear-Time Sequence Modeling with Selective State Spaces}. arXiv:2312.00752, 2023.
\bibitem{ttr} K.~A.~Wang, J.~Shi, E.~B.~Fox. \emph{Test-time regression: a unifying framework for designing sequence models with associative memory}. arXiv:2501.12352, 2025.
\bibitem{gdn} S.~Yang, J.~Kautz, A.~Hatamizadeh. \emph{Gated Delta Networks: Improving Mamba2 with Delta Rule}. arXiv:2412.06464, 2025.
\bibitem{trm} A.~Jolicoeur-Martineau. \emph{Less is More: Recursive Reasoning with Tiny Networks}. arXiv:2510.04871, 2025.
\bibitem{recdepth} J.~Geiping et al. \emph{Scaling up Test-Time Compute with Latent Reasoning: A Recurrent Depth Approach}. arXiv:2502.05171, 2025.
\bibitem{flesch} T.~Flesch, D.~G.~Nagy, A.~Saxe, C.~Summerfield. \emph{Modelling continual learning in humans with Hebbian context gating and exponentially decaying task signals}. arXiv:2203.11560, 2022.
\bibitem{dendrites} A.~Iyer et al. \emph{Avoiding Catastrophe: Active Dendrites Enable Multi-Task Learning in Dynamic Environments}. arXiv:2201.00042, 2022.
\bibitem{cgar} M.~Choraria et al. \emph{Context-Gated Associative Retrieval: From Theory to Transformers}. arXiv:2605.10970, 2026.
\bibitem{mom} J.~Du et al. \emph{MoM: Linear Sequence Modeling with Mixture-of-Memories}. arXiv:2502.13685, 2025.
\bibitem{raven} A.~Afzal, A.~Bick et al. \emph{Raven: High-Recall Sequence Modeling with Sparse Memory Routing}. arXiv:2607.25357, 2026.
\bibitem{routing} A.~Basu. \emph{When Does Content-Based Routing Work? Representation Requirements for Selective Attention in Hybrid Sequence Models}. arXiv:2603.20997, 2026.
\bibitem{hmrnn} J.~Chung, S.~Ahn, Y.~Bengio. \emph{Hierarchical Multiscale Recurrent Neural Networks}. arXiv:1609.01704, 2017.
\bibitem{nested} A.~Behrouz et al. \emph{Nested Learning: The Illusion of Deep Learning Architectures}. arXiv:2512.24695, 2025.
\bibitem{switch} W.~Fedus, B.~Zoph, N.~Shazeer. \emph{Switch Transformers}. JMLR 23(120), 2022.
\bibitem{zoology} S.~Arora et al. \emph{Zoology: Measuring and Improving Recall in Efficient Language Models}. arXiv:2312.04927, 2023.
\bibitem{bietti} A.~Bietti, V.~Cabannes, D.~Bouchacourt, H.~J\'egou, L.~Bottou. \emph{Birth of a Transformer: A Memory Viewpoint}. arXiv:2306.00802, 2023.
\bibitem{edelman} B.~L.~Edelman, E.~Edelman, S.~Goel, E.~Malach, N.~Tsilivis. \emph{The Evolution of Statistical Induction Heads: In-Context Learning Markov Chains}. arXiv:2402.11004, 2024.
\end{thebibliography}

\end{document}
